% =============================================================================
%  一般线性群的局部 Langlands 对应 —— 中译注本
%
%  原文：Zijian Yao, "Notes on Local Langlands for GLn",
%        Harvard 学习讨论班讲义（2017 年秋），36 页；作者自述为极粗糙的笔记。
%  编译：xelatex main.tex 三遍
%
%  体例：
%    * 底本的分节编号 1.1.1、2.3.1 等照录，定理/引理/推论共用一套按节重排的编号；
%    * 底本未完成的节（1.1.8、1.1.9、3.3.2、3.4、3.6.2 等）照实标注、不代拟；
%    * 「译注」（橙框）与「本章要点」（蓝框）为译者所加，非原文内容；
%    * 原注以上标数字标记，集中照录于文末「原注」。
% =============================================================================
\documentclass[11pt,a4paper]{article}

%% 导言区 —— SY02-43《一般线性群的局部 Langlands 对应》中译注本
%% 底本：Zijian Yao, "Notes on Local Langlands for GLn",
%%       Harvard 学习讨论班讲义（2017 年秋），36 页，pdfTeX 排版。
%% 编译：xelatex main.tex（三遍，待目录稳定）
\usepackage[UTF8, fontset=windows, heading=true]{ctex}
\usepackage{amsmath}
\usepackage{mathtools}
\usepackage{amssymb}
\usepackage{amsthm}
\usepackage{bm}
\usepackage[a4paper,left=2.6cm,right=2.6cm,top=2.6cm,bottom=2.6cm]{geometry}
\usepackage{graphicx}
\usepackage{xcolor}
\usepackage[most]{tcolorbox}
\usepackage{enumitem}
\usepackage{booktabs}
\usepackage{array}
\usepackage[font=small,labelfont=bf,labelsep=quad]{caption}
\usepackage[colorlinks=true,
            linkcolor=blue!45!black,
            citecolor=blue!45!black,
            urlcolor=blue!45!black]{hyperref}

\setlength{\emergencystretch}{3em}
\allowdisplaybreaks[3]
\linespread{1.12}
\renewcommand{\figurename}{图}
\renewcommand{\tablename}{表}
\ctexset{
  section/name       = {},
  section/number     = \arabic{section},
  subsection/name    = {},
  subsection/number  = \thesection.\arabic{subsection},
  subsubsection/name = {},
}

%% ---------- 定理族：底本中 定理/引理/推论/命题 共用一个按节重排的计数器 ----------
\theoremstyle{plain}
\newtheorem{nbthm}{定理}[section]
\newtheorem{theorem}[nbthm]{定理}
\newtheorem{lemma}[nbthm]{引理}
\newtheorem{corollary}[nbthm]{推论}
\newtheorem{proposition}[nbthm]{命题}
\newtheorem*{theorem*}{定理}

\theoremstyle{definition}
\newtheorem*{definition}{定义}
\newtheorem*{example}{例}
\newtheorem*{remark}{注}
\newtheorem*{corollarydefinition}{推论-定义}

\renewcommand{\proofname}{证明}

%% ---------- 「本章要点」（译者所加） ----------
\newtcolorbox{keytakeaways}[1][本章要点]{
  enhanced,
  colback  = blue!3!white,
  colframe = blue!35!black,
  boxrule  = 0.4pt,
  arc      = 2pt,
  left     = 7pt, right = 7pt, top = 5pt, bottom = 5pt,
  title    = {#1\,{\normalfont\scriptsize（译者所加）}},
  fonttitle= \bfseries\small,
  before skip = 10pt, after skip = 10pt,
}

%% ---------- 译注（译者所加） ----------
\newif\ifshowyizhu
\showyizhutrue
\newcounter{yizhu}
\makeatletter
\newenvironment{yizhu}[1][]{%
  \ifshowyizhu
    \refstepcounter{yizhu}%
    \par\addvspace{5pt}\noindent
    \begin{tcolorbox}[enhanced, breakable,
      colback  = orange!4!white,
      colframe = orange!55!black,
      boxrule  = 0.3pt,
      arc      = 1.5pt,
      left     = 6pt, right = 6pt, top = 4pt, bottom = 4pt,
      title    = {\small\bfseries 译注~\theyizhu
                  \if\relax\detokenize{#1}\relax\else\ ——#1\fi}]%
    \small
  \else
    \setbox\z@\vbox\bgroup\small
  \fi
}{%
  \ifshowyizhu
    \end{tcolorbox}%
  \else
    \egroup
  \fi
}
\makeatother

%% 底本用上标数字标脚注，译文以 \en{n} 复现并集中于文末「原注」
\newcommand{\en}[1]{\textsuperscript{#1}}

%% 记号宏 —— SY02-43
%% 原则：宏名只用于高频、易错记号；正文数学仍以标准 LaTeX 书写，便于与原文对照。
\newcommand{\bbA}{\mathbb{A}}
\newcommand{\bbC}{\mathbb{C}}
\newcommand{\bbF}{\mathbb{F}}
\newcommand{\bbN}{\mathbb{N}}
\newcommand{\bbP}{\mathbb{P}}
\newcommand{\bbQ}{\mathbb{Q}}
\newcommand{\bbR}{\mathbb{R}}
\newcommand{\bbZ}{\mathbb{Z}}
\newcommand{\Qbar}{\overline{\bbQ}}
\newcommand{\Ql}{\bbQ_\ell}
\newcommand{\Fq}{\bbF_q}
\newcommand{\Gm}{\mathbb{G}_m}
\newcommand{\Ga}{\mathbb{G}_a}

\newcommand{\GL}{\mathrm{GL}}
\newcommand{\SL}{\mathrm{SL}}
\newcommand{\sO}{\mathrm{O}}
\newcommand{\SO}{\mathrm{SO}}
\newcommand{\Sp}{\mathrm{Sp}}
\newcommand{\sT}{\mathrm{T}}
\newcommand{\sB}{\mathrm{B}}
\newcommand{\sU}{\mathrm{U}}
\newcommand{\sN}{\mathrm{N}}
\newcommand{\sW}{\mathrm{W}}
\newcommand{\sI}{\mathrm{I}}
\newcommand{\sP}{\mathrm{P}}
\newcommand{\Rep}{\mathrm{Rep}}

\newcommand{\Hom}{\mathrm{Hom}}
\newcommand{\End}{\mathrm{End}}
\newcommand{\Aut}{\mathrm{Aut}}
\newcommand{\Ind}{\mathrm{Ind}}
\newcommand{\ind}{\mathrm{ind}}
\newcommand{\Res}{\mathrm{Res}}
\newcommand{\tr}{\mathrm{tr}}
\newcommand{\rk}{\mathrm{rank}}
\newcommand{\diag}{\mathrm{diag}}
\newcommand{\Spec}{\mathrm{Spec}}
\newcommand{\Frob}{\mathrm{Frob}}
\newcommand{\Art}{\mathrm{Art}}
\newcommand{\Gal}{\mathrm{Gal}}
\newcommand{\Tr}{\mathrm{Tr}}
\newcommand{\Ker}{\ker}
\newcommand{\supp}{\mathrm{supp}}
\newcommand{\ad}{\mathrm{ad}}
\newcommand{\St}{\mathrm{St}}
\newcommand{\triv}{\mathrm{triv}}
\newcommand{\sgn}{\mathrm{sgn}}
\newcommand{\e}{\mathrm{e}}
\newcommand{\vol}{\operatorname{Vol}}

\newcommand{\sWD}{\mathrm{wd}}
\newcommand{\WD}{\mathrm{WD}}
\newcommand{\Asc}{\mathrm{Asc}}
\newcommand{\Girr}{G^{\mathrm{irr}}}
\newcommand{\An}{A_n}
\newcommand{\Repinf}{\Rep^{\infty}}
\newcommand{\Repadm}{\Rep^{\mathrm{adm}}}

\newcommand{\wt}[1]{\widetilde{#1}}
\newcommand{\ol}[1]{\overline{#1}}
\newcommand{\wh}[1]{\widehat{#1}}
\newcommand{\abs}[1]{\left|#1\right|}
\newcommand{\norm}[1]{\left\|#1\right\|}
\newcommand{\slt}{\mathfrak{sl}}
\newcommand{\gl}{\mathfrak{gl}}
\newcommand{\Lf}[1]{L(#1,s)}
\newcommand{\twist}{\otimes}

%% 底本高频记号
\newcommand{\WK}{W_K}
\newcommand{\GK}{G_K}
\newcommand{\IK}{I_K}
\newcommand{\PK}{P_K}
\newcommand{\evareps}{\epsilon}
\newcommand{\rec}{\mathrm{rec}}


\begin{document}

% =============================================================================
%  扉页、摘要、译者说明
% =============================================================================
\begin{titlepage}
\centering
\vspace*{2.6cm}
{\Huge\bfseries 一般线性群的\\[0.25em] 局部 Langlands 对应\par}
\vspace{1.0em}
{\large\itshape Notes on Local Langlands for $\GL_n$\par}
\vspace{2.2em}
{\large Zijian Yao\par}
\vspace{0.4em}
{\normalsize 哈佛大学学习讨论班讲义，2017 年秋\par}
\vspace{2.6em}
{\large 中译注本\par}
\vspace{0.6em}
{\normalsize 数译 · Langlands 导读系列 · SY02-43\par}
\vfill
{\small
\begin{tabular}{@{}rl@{}}
底本 & Zijian Yao, \emph{Notes on Local Langlands for $\GL_n$},\\
     & 哈佛学习讨论班讲义（2017 年秋），36 页（pdfTeX 排版）\\
译者 & 数译（DSH agent：lead）\\
版本 & v1.0，2026 年 11 月 27 日\\
体例 & 分节编号 1.1.1、2.3.1 等照底本；\\
     & 定理、引理、推论共用一套按节重排的编号；\\
     & 底本未写完的小节照实标注、不代拟；\\
     & 译注（橙色框）与「本章要点」（蓝色框）为译者所加\\
\end{tabular}}
\end{titlepage}

\section*{摘要}
\addcontentsline{toc}{section}{摘要}

这是 2017 年秋季哈佛一次学习讨论班的极粗糙笔记，由作者最初几次报告扩充而来。
由于报告人与听众中都没有专家，笔记里多半有我们尚未察觉的错误，阅读时请谨慎；
若发现关键性的误解，作者希望读者以电子邮件告知。

全书分六节。第 1 节整理 Weil--Deligne 表示与 $l$-进表示，以及两者之间的
$\mathrm{WD}$ 函子。第 2 节从光滑表示、容许表示、Hecke 模讲到超尖点表示与
Jacquet 函子，最后陈述 $\GL_n$ 的局部 Langlands 对应。第 3 节照 Bernstein--Zelevinsky
的分类，把 $\GL_n$ 的容许表示组织成 Zelevinsky 段的乘积，并讨论平方可积、
缓增、一般（generic）表示与 $K$-球面表示。第 4 节讨论惯性 Langlands 对应与
Bernstein 分解、类型论。第 5、6 节在底本中只有标题，正文告缺。

\begin{yizhu}[关于「粗糙笔记」的说明]
底本扉页自述 these notes are extremely rough，且多处明确写着
「to be completed in second draft」。译文对这类空缺一律照实保留（标注
\emph{底本此处未写}），不替作者补写内容；凡译者补出的背景，一律放进橙色译注框。
\end{yizhu}

\section*{译者说明}
\addcontentsline{toc}{section}{译者说明}

\paragraph{底本与体例}
底本为 Zijian Yao 2017 年秋在哈佛学习讨论班的讲义
\emph{Notes on Local Langlands for $\GL_n$}，36 页，无正式出版信息，
作者自印 pdfTeX 稿（文件内嵌日期 2019-10-11）。底本的特征有三点需要在译文里交代：

\begin{enumerate}[leftmargin=2em]
\item \textbf{分层写作}。作者按「第一稿写骨架、第二稿补内容」推进，因此第 1、2 节基本完整，
      第 3 节已开始出现空小节（3.3.2、3.4、3.6.2），第 5、6 节只剩标题。
      译文保留全部小节编号与标题，空缺处标注\emph{底本此处未写}。
\item \textbf{大量「注」「例」不编号}。底本的定义、注、例一律不编号，
      编号只给定理、引理、推论、命题，且四者共用一个按节重排的计数器
      （例如 1.3 是定理，1.4 是推论）。译文照此办理。
\item \textbf{记号不统一处照录}。底本在 1.1.4 用 $V(-1)$ 表示形式符号意义下的扭转，
      又用 $\epsilon: G_K \to L^\times$ 表示对应的非分歧特征；正文两处都保留，
      但在译注里点明这是「形式符号 + 具体实现」两层意思。
\end{enumerate}

\paragraph{术语}
译文术语以本库已冻结的译名为准，并与香蕉空间（bananaspace.org）现有词条交叉核对：
「局部 Langlands 对应」「惯性群」「尖点表示」「导子」「非分歧」「分歧」「驯分歧」
在该站有页面或导航框用例，直接沿用；「超尖点表示」「容许表示」「Weil--Deligne 表示」
在该站无专条，沿本库既有译名（超尖点、容许的、Weil--Deligne 表示）。
逐条对照见 \texttt{资料卡/04-术语文献勘误.md}。

\paragraph{译注的分工}
橙色框只做三件事：补底本省略的推理、点明记号冲突、交代该概念在中文文献里的别名。
凡涉及数学内容的判断，都以「底本如此」与「译者按」两类措辞区分，
不与原文陈述混同。

\paragraph{致谢与免责}
底本作者明言笔记可能有错。译文对底本中明显的笔误（如把 $\GL_n$ 写成 $\GL_2$、
把 $\sigma$ 写成 $\psi$）在译注中列出，但不改动正文公式；读者若发现译文自身的问题，
请以底本为准。


\newpage
\tableofcontents
\newpage

% =============================================================================
%  §1 Weil--Deligne 表示（底本 p.2--9）
% =============================================================================
\section{Weil--Deligne 表示}

本节中固定一个非 Archimedes 局部域 $K$，其剩余特征为 $p$。记 $k$ 为其剩余域，
$q=\#k=p^{f}$。固定一个素元 $\varpi_K\in\mathcal{O}_K$。最后，设
$G_K=\Gal(K^{s}/K)$ 为固定的可分闭包 $K^{s}$ 在 $K$ 上的 Galois 群。

\subsection{Weil--Deligne 表示}

\subsubsection{Weil 群 $W_K$}

按定义我们有交换图
\[
\begin{array}{ccccccccc}
0 & \longrightarrow & I_K & \longrightarrow & W_K & \longrightarrow & \Frob^{\mathbb{Z}} & \longrightarrow & 0\\[2pt]
 & & \Big\downarrow\simeq & & \Big\downarrow & & \Big\downarrow\simeq & & \\[2pt]
0 & \longrightarrow & I_K & \longrightarrow & G_K & \longrightarrow & G_k & \longrightarrow & 0 .
\end{array}
\]
这里 $W_K$ 的拓扑取成使 $I_K\subset W_K$ 为开子群、并带上它通常的拓扑，而
$\Frob^{\mathbb{Z}}\simeq\mathbb{Z}$ 取离散拓扑。特别地，这\emph{不是}从 $G_K$ 诱导的子空间拓扑。

\begin{remark}
一般地，若 $K$ 是局部域或整体域，与 $K$ 相伴的 $W_K$ 的定义涉及一个三元组
$(W_K,\iota,\{r_E\})$：其中 $W_K$ 是拓扑群，$\iota\colon W_K\to G_K$ 是像稠密的连续同态，
而对每个有限扩张 $E/F$ 有 $r_E\colon C_E\xrightarrow{\ \sim\ }W^{\mathrm{ab}}_E$，
这里 $C_E$ 在局部情形等于 $E^{\times}$、在整体情形等于 $E^{\times}\backslash\mathbb{A}^{\times}_E$；
这些数据满足若干自然条件。可以证明这样的 $W_K$ 存在且唯一。在非 Archimedes 局部情形，
它与上面的定义一致：$\iota$ 是自然的包含，$r_E$ 是局部类域论给出的同构。
\end{remark}

\subsubsection{驯特征标}

固定 $K^{s}$ 中一族相容的本原单位根 $\{\zeta_n\}_{p\nmid n}$。于是有 $K$ 上的
最大非分歧扩张 $K^{\mathrm{nr}}=\bigcup_{p\nmid n}K(\zeta_n)$ 与最大驯分歧扩张
$K^{t}=\bigcup_{p\nmid n}K^{\mathrm{nr}}(\varpi_K^{1/n})$。野惯性群
$P_K=G(K^{s}/K^{t})$ 是 $I_K$ 中唯一的 pro-$p$ Sylow 子群，而\emph{驯特征标}
（依赖于 $\{\zeta_n\}$ 的选取）是如下同构：
\[
t\colon I_K/P_K\xrightarrow{\ \sim\ }\prod_{l\neq p}\mathbb{Z}_l,
\qquad
\sigma\longmapsto t(\sigma)\ \mathrm{mod}\ n
\quad\text{其中}\quad
\frac{\sigma(\varpi_K^{1/n})}{\varpi_K^{1/n}}=\zeta_n^{t(\sigma)} .
\]
此外，对 $l\neq p$，令 $K^{t,l}=K^{\mathrm{nr}}(\varpi_K^{1/l^{\infty}})$ 为最大驯分歧扩张的
$l$-部分，则有（作为 Abel 群的同构）
\[
t_l\colon I_{K,l}/P_K=G(K^{t,l}/K^{\mathrm{nr}})\xrightarrow{\ \sim\ }\mathbb{Z}_l ,
\]
它与驯特征标到 $\mathbb{Z}_l$ 的投影一致。

$G_K$（或 $W_K$）通过共轭作用在 $I_K/P_K$ 上（因为 $P_K$ 是唯一的 pro-$p$ Sylow 子群）。
在每个 $\mathbb{Z}_l$ 分量上，该作用由 $l$-进分圆特征标 $\chi=\chi_l$ 给出：
\[
I_{K,l}/P_K\xrightarrow{\ \sim\ }\mathbb{Z}_l(1),
\]
其中的 $G_K$-作用是共轭。换言之，对任意 $\sigma\in G_K$、$\tau\in I_{K,l}$，有
\[
t_l(\sigma\tau\sigma^{-1})=\chi(\sigma)\,t_l(\tau).
\]

\begin{yizhu}[$\mathbb{Z}_l(1)$ 的读法]
底本把 $I_{K,l}/P_K\cong\mathbb{Z}_l$ 与 $\mathbb{Z}_l(1)$ 等同，靠的正是「共轭作用 =
乘 $\chi_l$」。因此这里的 $\mathbb{Z}_l(1)$ 不是抽象群，而是带 $G_K$-作用的模：
$G_K$ 通过 $\chi_l$ 作用在 $\mathbb{Z}_l$ 上。
\end{yizhu}

\subsubsection{Weil--Deligne 表示}

设 $L$ 是特征 $0$ 的域，$V$ 是 $L$ 上的有限维向量空间。

\begin{definition}
$W_K$ 在 $V$ 上的一个\emph{Weil--Deligne 表示}是一对 $(\rho,N)$，其中
\begin{enumerate}[label=(\arabic*),leftmargin=2.2em]
\item $\rho$ 是 $W_K$ 在 $V$ 上的连续表示（$L$ 取离散拓扑），等价地，$I_K$ 的像有限；
\item $N\in\End(V)$ 是自同态，使得对任意 $\sigma\in W_K$ 有
\[
\sigma N\sigma^{-1}=q^{-v_K(\sigma)}N ,
\]
其中 $q=\#k$，而 $v_K$ 由 $\sigma\mapsto\Frob_K^{v_K(\sigma)}\in W_K/I_K
\xrightarrow{\ \sim\ }\Frob^{\mathbb{Z}}_K$ 确定。
\end{enumerate}
\end{definition}

\begin{definition}
记 $\mathrm{wd}=\mathrm{wd}_L(W_K)$ 为 Weil--Deligne 表示构成的范畴，其态射是
与自同态 $N$ 相容的表示态射。
\end{definition}

\subsubsection{基本性质}

设 $(\rho,N)\colon W_K\to\GL(V)$ 是如上给出的 Weil--Deligne 表示。特别地，对任一
Frobenius 的提升 $\varphi\in W_K$，有
\[
\varphi N\varphi^{-1}=q^{-1}N .
\]
我们记录如下性质：

\begin{enumerate}[label=(\arabic*),leftmargin=2.2em]
\item $N$ 必然是幂零的：特征多项式满足
\[
f(X)=\det(I_n\cdot X-N)=\det(I_n\cdot X-q^{-1}N)=q^{-n}f(qX).
\]
于是 $f(X)=X^{n}$，这迫使 $N$ 幂零。

\item $N$ 显然不是 Galois 等变的；事实上 $N$ 诱导出一个 Galois 等变的映射
\[
N(-1)\colon V\longrightarrow V(-1),
\]
其中 $V(-1)$（作为一个形式符号）表示 $V$ 用非分歧特征标
$\epsilon\colon G_K\to L^{\times}$、$\varphi\mapsto q$ 所作的扭转，$\varphi$ 是 Frobenius 的一个提升。
若 $L/\mathbb{Q}_l$ 且 $l\neq p$，则 $(-1)$ 与 Tate 扭转的通常记法一致。

\item 设 $r\colon W_K\to\GL(V)$ 是 $W_K$ 的连续表示（$V$ 上取离散拓扑），则定义
$W_K$ 在 $V^{m}$ 上的 Weil--Deligne 表示 $\mathrm{Sp}_m(r)$ 为
\[
\mathrm{Sp}_m(r):=r\oplus r(1)\oplus\cdots\oplus r(m-1),
\]
其中 $r(i)$ 是 $r$ 用 $\chi^{-1}$ 所作的扭转；再定义 $N\in\End(V^{m})$ 使其对
$i\le m-2$ 把 $r(i)$ 同构地送到 $r(i+1)$，而在 $r(m-1)$ 上取 $N=0$。等价地，
\[
N(-1)\colon \mathrm{Sp}_m(r)=r\oplus\cdots\oplus r(m-1)\longrightarrow
\mathrm{Sp}_m(r)(-1)=r(-1)\oplus\cdots\oplus r(m-2)
\]
在 $r\oplus\cdots\oplus r(m-2)$ 上是恒等、在 $r(m-1)$ 上是 $0$。

\item 对使 $N^{m}=0$ 的最小 $m$ 作归纳，可以证明：每个使 $\rho$ 半单的
Weil--Deligne 表示 $(\rho,N)$ 都同构于若干形如 $\mathrm{Sp}_m(r)$ 的 Weil--Deligne 表示的直和，
其中 $r$ 是 $W_K$ 的不可约表示、且 $r(I_K)$ 有限。
\end{enumerate}

\begin{yizhu}[关于 $r(i)$ 的扭转方向]
底本写作「$r(i)$ 是 $r$ 用 $\chi^{-1}$ 所作的扭转」，而第 (2) 条又把 $V(-1)$ 定义为
用 $\epsilon(\varphi)=q$ 扭转。两处方向相反，读者容易在此处错一次符号。
译者按：本库统一采用「$V(1)=V\otimes\chi$，$V(-1)=V\otimes\chi^{-1}$」的 Tate 约定，
于是 $r(i)=r\otimes\chi^{-i}$ 与 $\mathrm{Sp}_m(r)$ 中 $N(-1)$ 把 $r(i)$ 送到 $r(i+1)$ 相容；
底本第 (2) 条的 $V(-1)$ 应读作「$N$ 把权 $-1$ 的 Tate 扭转送回权 $0$」这一形式陈述，
其符号以第 (3) 条的显式描述为准。
\end{yizhu}

\subsubsection{关于半单表示的注记}

本段讨论中设 $V$ 是任意域 $k$ 上的向量空间，$H<G$ 是子群，$\rho\colon G\to\GL(V)$ 是
$G$ 的表示。我们想知道 $\rho$ 的半单性与 $\rho|_H$ 的半单性之间的关系。

\begin{lemma}
\leavevmode
\begin{enumerate}[label=(\arabic*),leftmargin=2.2em]
\item 若 $V$ 有限维、$H\triangleleft G$ 是正规子群，则 $\rho$ 半单蕴含 $\rho|_H$ 半单。
\item 若 $H<G$ 的指数为 $r$，且 $r$ 在系数域 $k$ 中可逆，则 $\rho$ 半单当且仅当 $\rho|_H$ 半单。
\end{enumerate}
\end{lemma}

\begin{proof}
(1) 不妨设 $\rho$ 不可约。取 $W=\bigoplus_i W_i$，其中 $\{W_i\}$ 是 $V$ 的互不相同的
$H$-子表示，则由 $H$ 的正规性 $W$ 是 $V$ 的 $G$-子表示。由 $V$ 有限维知 $W\neq 0$，
于是 $W=V$，即 $\rho|_H=\bigoplus W_i$ 半单。

(2) 设 $G=\bigsqcup_{i=1,\dots,r}Hg_i$。

(a) 先设 $\rho$ 不可约。任一 $v\neq 0$ 作为单 $k[G]$-模生成 $V$，故 $\{g_iv\}$ 作为
$k[H]$-模生成 $V$。现取 $W\subset V$ 为极大的 $k[H]$-子模（因 $V$ 有限生成而存在）。
则每个 $V/g_iW$ 是单 $k[H]$-模，且有自然的 $k[H]$-线性映射 $V\to\bigoplus_i V/g_iW$，
其核是 $V$ 的真 $k[G]$-子模，故为 $0$。于是 $V\hookrightarrow\bigoplus_i V/g_iW$ 半单。

(b) 「若」的方向是 Maschke 定理。设 $\rho|_H$ 半单，$W\subset V$ 是 $G$-子表示，
从而也是半单的 $H$-子表示。取 $H$-线性投影 $\pi\colon V\to W$，对 $\pi$ 取平均得
$G$-线性投影
\[
\pi^{G}\colon V\longrightarrow W,\qquad
v\longmapsto\frac{1}{r}\sum_{i=1,\dots,r}g_i^{-1}\pi(g_i\cdot v),
\]
故得证。（要点在于：即使 $g_i$ 与 $h$ 不可交换，$g_i^{-1}\circ\pi\circ g_i$ 仍与 $h$ 交换，
因此 $\pi^{G}$ 是 $H$-线性的，从而是 $G$-线性的。）
\end{proof}

\subsubsection{$\varphi$-半单表示}

\begin{lemma}
设 $\rho$ 是 $W_K$ 在 $V$ 上的表示，且 $\rho(I_K)$ 有限。为简单起见设 $L$ 代数闭。则下列等价：
\begin{enumerate}[label=(\arabic*),leftmargin=2.2em]
\item 对一切 $\sigma\in W_K$，$\rho(\sigma)$ 半单；
\item 对 Frobenius 的（任一）提升 $\varphi$，$\rho(\varphi)$ 半单；
\item $\rho$ 是半单表示。
\end{enumerate}
\end{lemma}

\begin{proof}[略证]
注意 $\operatorname{char}L=0$。由于 $\operatorname{im}\rho(I_K)$ 有限，我们断言：对任意
$\sigma\in W_K$，存在 $m\in\mathbb{Z}_{>0}$ 使 $\rho(\sigma^{m})$ 落在 $\rho(W_K)$ 的中心里。
为证此断言，取 $\tau_1,\dots,\tau_r\in I_K$ 生成 $\rho(I_K)$，并令 $\tau_0=\varphi$ 是
Frobenius 的提升。对每个 $\tau_i$，考虑 $\rho(\tau_i^{n}\sigma\tau_i^{-n})$（$n\in\mathbb{Z}_{>0}$），
由 $\rho(I_K)$ 的有限性知存在 $m_i$ 使 $\rho(\tau_i^{m_i})$ 与 $\rho(\sigma)$ 交换；
取 $m=m_0m_1\cdots m_r$ 即可。

于是对任意 $\sigma\in W_K$，存在 $n\ge 0$、$n'>0$ 使 $\rho(\sigma^{n'})=\rho(\varphi^{n})$。
由于 $L$ 特征为 $0$，线性算子 $T$ 的半单性等价于 $T$ 的幂的半单性（注意这在正特征下
一般不成立），故 (1) 与 (2) 等价。又 $\rho(\varphi)$ 生成 $\rho(W_K)$ 中的有限指数子群，
而 $L$ 特征为 $0$，由引理 1.1 第 (2) 部分，$\rho(W_K)$ 在 $V$ 上半单作用当且仅当
$\langle\rho(\varphi)\rangle$ 半单作用，当且仅当 $\rho(\varphi)$ 半单。
\end{proof}

\begin{remark}
注意上一条引理中我们假定了 $L$ 代数闭。这是因为按我们的约定，线性算子称为半单，
指的是它「潜在可对角化」，即在某次基变换后可对角化。对一般的 $L$，只需把第 (3) 条
中的说法换成「潜在半单」。
\end{remark}

\begin{definition}
Weil--Deligne 表示 $(\rho,N)\in\mathrm{wd}_K$ 称为 \emph{$\varphi$-半单的}
（Frobenius 半单的），若 $\rho$ 半单；等价地，若 $\rho(\varphi)$ 半单。
\end{definition}

\begin{remark}
\leavevmode
\begin{enumerate}[label=(\arabic*),leftmargin=2.2em]
\item 设 $(\rho,N)\in\mathrm{wd}$。若 $\rho$ 不可约则 $N$ 必为 $0$，因为 $\ker N$
在 $W_K$ 下稳定。
\item $\varphi$-半单表示一般\emph{不是}范畴 $\mathrm{wd}$ 中的半单对象。
$\mathrm{wd}$ 中的半单对象是单对象的直和，即那些满足 $N=0$ 的 $\varphi$-半单表示。
\end{enumerate}
\end{remark}

\begin{remark}[Frobenius 半单化]
设 $(\rho,N)$ 是 $W_K$ 的 Weil--Deligne 表示，$\rho^{\mathrm{ss}}$ 是 $\rho$ 的半单化，
则 $(\rho^{\mathrm{ss}},N)$ 仍是 Weil--Deligne 表示，称为 $(\rho,N)$ 的
\emph{$\varphi$-半单化}。
\end{remark}

\subsubsection{张量}

设 $(\rho_1,N_1)$ 与 $(\rho_2,N_2)$ 是两个 Weil--Deligne 表示，则它们的张量积是
\[
(\rho_1\otimes\rho_2,\ N_1\otimes\mathrm{id}+\mathrm{id}\otimes N_2).
\]
类似地，$\Hom(\rho_1,\rho_2)=(r,N)$ 由下式给出：
\[
\bigl(r(\tau)\varphi\bigr)v_1=\rho_2(\tau)\cdot\varphi\bigl(\rho_1(\tau)^{-1}v_1\bigr),
\qquad
(N\varphi)v_1=N_2\cdot\varphi(v_1)-\varphi(N_1\cdot v_1).
\]

\subsubsection{Weil--Deligne 群概形}

把 $W_K$ 看作 $\mathbb{Q}$ 上的（常）群概形，它取得自极限
$\varprojlim_{I'\triangleleft I_K}W_K/I'$，其中 $I'$ 遍历 $I_K$ 的所有开正规子群。
这是 $\mathbb{Q}$ 上的无限型群概形。令
\[
\widehat{W}_K:=\mathbb{G}_a\rtimes W_K ,
\]
其中 $W_K$ 在 $\mathbb{G}_a$ 上的作用由 $wxw^{-1}=|w|_K x$ 给出。

\begin{center}
\emph{底本此处未写}（作者原注：to be completed in second draft）。
\end{center}

\subsubsection{Jacobson--Morozov}

\begin{center}
\emph{底本此处未写}（作者原注：to be completed in second draft）。
\end{center}

\subsubsection{Tate 与 Deligne--Langlands 的 $\epsilon$-因子}

我们需要为 Weil--Deligne 表示 $(\rho,N)$ 定义 $L$-导子与 $\epsilon$-导子，
在此过程中还要定义导子 $c((\rho,N))$。先对 $W_K$ 的表示做这件事。回忆：对我们而言，
$W_K$ 的表示指的是关于目标上离散拓扑连续的表示。

若固定一个同构 $L=\mathbb{Q}_l\xrightarrow{\ \sim\ }\mathbb{C}$，这些表示对应于
$W_K$ 的 Artin 表示。固定 $\varpi=\varpi_K$。

先回忆（Tate 论文中的结论）：对拟特征标 $\chi\colon K^{\times}\to\mathbb{C}$，Tate 定义局部
$L$-因子
\[
L(\chi)=
\begin{cases}
(1-\chi(\varpi))^{-1}, & \chi\ \text{非分歧},\\
1, & \chi\ \text{分歧}.
\end{cases}
\]
对 $\omega_s=\|\cdot\|^{s}$，$s\in\mathbb{C}$，$L(\chi\omega_s)$ 关于 $s$ 是亚纯的。
设 $\psi$ 是 $K$ 上非平凡加性特征标，$dx$ 是 $K$ 上的 Haar 测度。记 $d^{\times}x$ 为
$K^{\times}$ 上由 $\|x\|^{-1}dx$ 给出的 Haar 测度。最后，令
$\chi^{\vee}:=\chi^{-1}\omega_1=\chi^{-1}\|\cdot\|$ 为 $\chi$ 的扭对偶。
Tate 考虑 Fourier 变换 $\hat f(y)=\int f(x)\psi(xy)\,dx$，并证明
\[
\frac{\int\hat f^{\chi^{\vee}}(x)\,d^{\times}x}{L(\chi^{\vee})}
=\epsilon(\chi,\psi,dx)\cdot
\frac{\int f^{\chi}(x)\,d^{\times}x}{L(\chi)}
\]
其中 $\epsilon(\chi,\psi,dx)$ 是 $\chi$ 的非零全纯函数，称为 $\epsilon$-因子，
它依赖于 $\psi$ 与 $dx$。Tate 还证明：若令 $n(\psi)$ 为使 $\psi(\varpi^{-n}\mathcal{O}_K)=1$
的最大整数，$c(\chi)$ 为 $\chi$ 的导子，并取 $c\in K^{\times}$ 使 $v(c)=n(\psi)+c(\chi)$，则
\[
\epsilon(\chi,\psi,dx)=
\begin{cases}
\|c\|^{-1}\chi(c)\cdot\vol_{dx}(\mathcal{O}_K), & \chi\ \text{非分歧},\\[4pt]
\displaystyle\int_{c^{-1}\mathcal{O}_K^{\times}}\chi^{-1}(x)\psi(x)\,dx, & \chi\ \text{分歧}.
\end{cases}
\]
也就是说，他对特征标的 $\epsilon$-因子给出了显式公式。

对高维表示，（据 [10]）只有一条属于 Langlands 的存在性定理（Deligne 给了简短证明）。

\begin{theorem}
存在唯一的函数 $\epsilon$，对每个局部域 $K$、$W_K$ 的表示 $\rho$、非平凡加性特征标 $\psi$
与 $K$ 的 Haar 测度 $dx$，给出数 $\epsilon(\rho,\psi,dx)\in\mathbb{C}^{\times}$，使得
\begin{enumerate}[label=(\arabic*),leftmargin=2.2em]
\item 对特征标 $\chi$，$\epsilon(\chi,\psi,dx)$ 与上面的定义一致；
\item $\epsilon(\cdot,\psi,dx)$ 在 $W_K$ 表示的短正合列上是乘性的，特别地下降为
      Grothendieck 群上的映射 $\epsilon(\cdot,\psi,dx)\colon R(W_K):=K_0(\Rep W_K)\to\mathbb{C}^{\times}$；
\item $\epsilon$ 在下述意义下于 $0$ 次部分是归纳的：对每个域扩张 $E/L/K$ 与加性
      Haar 测度 $\nu_E$、$\mu_L$，有
\[
\epsilon(\Ind_{E/L}[r],\psi\circ\Tr_{L/K},\mu_L)
=\epsilon([r],\psi\circ\Tr_{E/K},\mu_E)
\]
对一切 $0$ 次元 $\ [r]\in R(W_K)$ 成立。
\end{enumerate}
\end{theorem}

\begin{remark}
设 $c(\rho)$ 是 $W_K$ 的表示 $(\rho,V)$ 的 Artin 导子，按定义等于
\[
c(\rho):=c(\chi_{\rho})=\sum_{i\ge 0}\frac{1}{[G_0:G_i]}\bigl(\dim V-\dim V^{G_i}\bigr),
\]
其中 $\chi_{\rho}$ 是该表示的特征标，$G_i$ 是 $I_K$ 的第 $i$ 个分歧子群。则可以证明
\[
\epsilon(\rho\cdot\omega_s,\psi,dx)=(\#k)^{-\bigl(c(\rho)+n(\psi)\dim V\bigr)\cdot s}\cdot
\epsilon(\rho,\psi,dx).
\]
\end{remark}

\subsubsection{导子、$L$-因子与 $\epsilon$-因子}

现在可以定义 Weil--Deligne 表示对应的不变量了。

\begin{definition}
设 $\rho=(\rho',N)$ 是 $W_K$ 的 Weil--Deligne 表示。设 $\psi$ 是 $K$ 上的加性特征标、
$dx$ 是 $K$ 上的加性 Haar 测度，$q=\#k$ 是剩余域的大小。
\begin{enumerate}[label=(\arabic*),leftmargin=2.2em]
\item 定义导子 $c(\rho)=c(\rho',N)$ 为
\[
c(\rho):=c(\rho')+\dim V^{I_K}-\dim(\ker N)^{I_K}.
\]
\item 设 $\Phi$ 是（任一）几何 Frobenius。定义 $L$-因子为
\[
L(\rho,s):=\det\Bigl(1-q^{-s}\Phi\ \big|\ (\ker N)^{I_K}\Bigr)^{-1}.
\]
\item 定义 $\epsilon$-因子为（下面省略了 $dx$ 的记号）
\[
\epsilon(\rho,s,\psi):=q^{-\bigl(c(\rho)+n(\psi)\dim V\bigr)\cdot s}\cdot
\det\Bigl(-\Phi\ \big|\ V^{I_K}/(\ker N)^{I_K}\Bigr)\cdot\epsilon(\rho',\psi,dx).
\]
\end{enumerate}
\end{definition}

特别地，从定义中挑出如下平凡的推论。

\begin{corollary}
\[
\epsilon(\rho,s,\psi)=q^{-\bigl(c(\rho)+n(\psi)\dim V\bigr)\cdot s}\cdot\epsilon(\rho,0,\psi).
\]
\end{corollary}

\begin{remark}
$\epsilon(\rho,s,\psi)$ 的定义与 [10] 略有不同：按我们的术语，Tate 把
$\epsilon$-因子定义为 $q^{n(\psi)\dim V\cdot s}\cdot\epsilon(\rho,s,\psi)$，
这似乎与上一页的注记不相容。我们的约定与 [11] 一致。
\end{remark}

\subsection{$l$-进表示}

\subsubsection{$l$-进 Galois 表示}

现在设 $L/\mathbb{Q}_l$ 是有限扩张，$l\neq p$。设 $V$ 是 $L$ 上的有限维向量空间。
固定一族相容的单位根 $\{\zeta_n\}_{p\nmid n}$ 与上面定义的 $l$-进驯特征标
$t_l\colon I_K/P_K\to\mathbb{Z}_l$。回忆 Grothendieck 的准单幂幺定理：

\begin{lemma}
设 $\rho\colon G_K\to\GL(V)$ 是 $G_K$ 的 $l$-进表示（即 $\rho$ 关于 $V$ 上的 $l$-进拓扑连续），则
\begin{enumerate}[label=(\arabic*),leftmargin=2.2em]
\item 存在开子群 $I'_K\subset I_K$（对应某个有限扩张 $K'/K$）与唯一确定的幂零算子
      $N\in\End(V)$，使得对一切 $\sigma\in I'_K$ 有 $\rho(\sigma)=\exp\bigl(t_l(\sigma)N\bigr)$；
\item 此外，对任意 $\tau=\sigma\varphi^{n}\in W_K$（其中 $\sigma\in I_K$，$\varphi$ 是 Frobenius 的
      固定提升），有 $\rho(\tau)N\rho(\tau^{-1})=q^{-n}N$。
\end{enumerate}
\end{lemma}

\begin{proof}[略证]
由连续性，$G_K$ 稳定某个 $\mathcal{O}_L$-格 $\Lambda\subset V$，模 $l$ 约化得
$\rho\colon G_K\to\Aut(\Lambda/l)$。其核 $\ker\rho=G_{K''}\subset G_K$ 是有限指数子群，
即 $K''/K$ 有限。由于 $\GL(\Lambda)$ 模 $l$ 约化的核是 pro-$l$ 的，知
$\rho|_{I_{K''}}$ 通过 $t_l\colon I_{K''}\to\mathbb{Z}_l(1)$ 分解。取 $\sigma\in I_{K''}$ 使
$t_l(\sigma)$ 拓扑生成 $\mathbb{Z}_l(1)$（作为 Abel 群）。

注意 $G_{K''}$ 在 $t_l(I_{K''})$ 上的共轭作用就是分圆特征标，即
$t_l(\varphi\sigma\varphi^{-1})=\chi(\varphi)t_l(\sigma)$，这迫使 $\rho(\sigma)$ 的特征值都是单位根。
于是存在有限扩张 $K'/K$ 使 $\rho(I'_K)$ 幂幺地作用。必要时把 $\sigma$ 换成 $I_{K'}$ 中新的
拓扑生成元（仍记作 $\sigma$），再令 $N=\log\rho(\sigma)$，它由幂幺性定义良好。
\end{proof}

\begin{corollary}
设 $\rho\colon G_K\to\GL(V)$ 是 $G_K$ 的 $l$-进表示。若 $\rho$ 半单，则 $\rho(I_K)$ 有限。
\end{corollary}

\begin{proof}
由引理 1.1 第 (1) 部分，$\rho|_{I_K}$ 半单。再把 $\rho$ 限制到引理 1.5 证明中的
$I_{K'}\triangleleft I_K$，则 $\rho|_{I_{K'}}$ 半单且幂幺，故为恒等。
\end{proof}

\subsubsection{函子 $\mathrm{WD}_{\zeta,\varphi}$}

设 $r\colon W_K\to\GL_n(L)$ 是关于 $L$ 上离散拓扑的连续表示，其中 $L/\mathbb{Q}_p$。
则 $r(\sigma)$ 的特征值（在 $\overline{\mathbb{Q}}_l$ 中）对一切 $\sigma\in W_K$ 都是 $l$-进单位，
当且仅当它们对 $r(\varphi)$（或对任一 $\tau\notin I_K$ 的 $r(\tau)$）是 $l$-进单位，
当且仅当 $\det\rho(\varphi)\in\mathcal{O}_L^{\times}$ 且 $\varphi$ 的特征多项式落在 $\mathcal{O}_L[X]$ 中。
此时称 $r$ 是 \emph{$l$-整的}或 \emph{$l$-有界的}。

\begin{corollary}
固定如上 $\zeta=\{\zeta_n\}$ 与 $\varphi$。则存在范畴等价
\[
\mathrm{WD}=\mathrm{WD}_{\zeta,\varphi}\colon \Rep_L(G_K)\longrightarrow
\mathrm{wd}^{l\text{-}\mathrm{int}}_L(W_K),
\]
从 $L$-向量空间上 $G_K$ 的 $l$-进表示范畴，到 $L$-向量空间上 $W_K$ 的 $l$-整
Weil--Deligne 表示范畴。函子 $\mathrm{WD}$ 由
\[
\mathrm{WD}(\rho)=(\rho',N),\qquad
\rho'(\sigma\varphi^{n}):=\rho(\sigma\varphi^{n})\exp\bigl(-t_l(\sigma)N\bigr)
\]
给出。
\end{corollary}

\begin{remark}
\leavevmode
\begin{enumerate}[label=(\arabic*),leftmargin=2.2em]
\item $\mathrm{WD}_{\zeta',\varphi'}$ 与 $\mathrm{WD}_{\zeta,\varphi}$ 自然同构。
\item $\mathrm{wd}^{l\text{-}\mathrm{int}}_L(W_K)$ 中的 $\varphi$-半单表示，对应于 $\rho(\varphi)$
      为半单算子的 $l$-进表示 $\rho$。
\item 设 $\mathrm{WD}(\rho)=(\rho',N)$，则 $\rho$ 不可约当且仅当 $N=0$ 且 $\rho'$ 不可约。
\end{enumerate}
\end{remark}

由上面两个推论，我们有：

\begin{corollary}
沿用上面的记号，$\rho\colon G_K\to\GL(V)$ 对应 $\mathrm{WD}(\rho)=(\rho',N)$，则下列等价：
\begin{enumerate}[label=(\arabic*),leftmargin=2.2em]
\item $\rho$ 在 $L$ 的某次基变换后半单；
\item $\rho$ 是 $\varphi$-半单的（即 $\varphi$ 半单地作用）且 $\rho(I_K)$ 有限；
\item $\rho'$ 是 $\varphi$-半单的且 $N=0$。
\end{enumerate}
\end{corollary}

\begin{proof}
显然 (1)$\Leftrightarrow$(3)$\Rightarrow$(2)。最后，(2)$\Rightarrow$(3) 由「基本性质」中紧接
Weil--Deligne 表示定义之后的最后一个断言得到；或者也可以这样看出 (2)$\Rightarrow$(1)：
(2) 蕴含在把 $L$ 换成 $\overline{L}$ 后，$\langle\rho(\varphi)\rangle$ 是半单表示，
而 $\langle\rho(\varphi)\rangle<\rho(G_K)$ 是有限指数的，于是可以应用引理 1.1。
\end{proof}

\subsubsection{Weil--Deligne 惯性类型与 Galois 惯性类型}

\begin{definition}
仍设 $V$ 是有限维 $L$-向量空间。
\begin{enumerate}[label=(\arabic*),leftmargin=2.2em]
\item 一个\emph{Weil--Deligne 惯性类型}（有时也简称\emph{惯性类型}）是指表示
\[
\tau\colon I_K\longrightarrow\GL(V),
\]
使得存在 $\varphi$-半单的 Weil--Deligne 表示 $(\rho',N)\colon W_K\to\GL(V)$ 满足
$\tau=\rho'|_{I_K}$。特别地，$\tau$ 的像有限，且不携带关于单值算子 $N$ 的任何信息。
\item 一个\emph{Galois 惯性类型}是指连续的 $\varphi$-半单表示
\[
\xi\colon I_K\longrightarrow\GL(V),
\]
它可以延拓为连续的（$l$-进）表示 $\rho\colon W_K\to\GL(V)$。
\end{enumerate}
\end{definition}

特别地，Galois 惯性类型限制为 Weil--Deligne 惯性类型：取 $\rho$ 为 $\xi$ 的某个延拓，
令 $(\rho',N)=\mathrm{WD}(\rho)$，再取 $\rho'|_{I_K}$。我们下面用 Galois 惯性类型的分类
更直接地描述这件事。

两种惯性类型都可以看作如下数据。$n$ 维（$\varphi$-半单的）Weil--Deligne 表示都形如
$\bigoplus_i\mathrm{Sp}_{m_i}(r_i)$，其中 $r_i$ 是 $W_K$ 的 $n_i$ 维不可约表示，
$\sum_i n_im_i=n$。

\begin{lemma}
设 $(r,N)=\bigoplus_{i=1}^{d}\mathrm{Sp}_{m_i}(r_i)$ 与
$(r',N')=\bigoplus_{i=1}^{d'}\mathrm{Sp}_{m'_i}(r'_i)$ 是 $W_K$ 的两个 $n$ 维
Weil--Deligne 表示，如上（即所有 $r_i$、$r'_i$ 都不可约）。设 $\rho$、$\rho'$ 是
$\mathrm{WD}^{-1}$ 给出的相应 $l$-进表示。
\begin{itemize}[leftmargin=1.6em]
\item 令 $\tau=r|_{I_K}$、$\tau'=r'|_{I_K}$ 是两个 Weil--Deligne 惯性类型；
\item 令 $\xi=\rho|_{I_K}$、$\xi'=\rho'|_{I_K}$ 是两个相应的 Galois 惯性类型。
\end{itemize}
则
\begin{enumerate}[label=(\arabic*),leftmargin=2.2em]
\item 作为 $I_K$ 的表示，$\tau\simeq\tau'$，当且仅当把 $\{r_i\}$（带重复）与 $\{r'_i\}$
      相差非分歧特征标的扭转后两者一致；
\item $\xi\simeq\xi'$，当且仅当 $d=d'$，且适当重排指标后 $m_i=m'_i$、
      $r_i\simeq r'_i\otimes\eta_i$（$\eta_i$ 是 $W_K$ 的非分歧特征标）。
\end{enumerate}
\end{lemma}

\begin{proof}
由引理 4.3 立即得到。
\end{proof}

\begin{center}
\textbf{总结（对应关系）。}
\end{center}

在上述对应与相应的 $\varphi$-半单部分之下，
\[
\Rep_L(G_K)\ \longleftrightarrow\ \mathrm{wd}^{l\text{-}\mathrm{int}}_L(W_K),
\qquad
\Rep^{\varphi\text{-}\mathrm{ss}}_L(G_K)\ \longleftrightarrow\
\mathrm{wd}^{\varphi\text{-}\mathrm{ss},\,l\text{-}\mathrm{int}}_L(W_K),
\]
我们有
\[
\Rep^{\mathrm{ss}}_L(G_K)=\bigl\{\rho\in\Rep^{\varphi\text{-}\mathrm{ss}}_L(G_K):
\rho(I_K)<\infty\bigr\}
\ \longleftrightarrow\
\bigl\{(\rho',N)\in\mathrm{wd}^{\varphi\text{-}\mathrm{ss},\,l\text{-}\mathrm{int}}_L(W_K):N=0\bigr\}.
\]

\begin{keytakeaways}
\begin{itemize}[leftmargin=1.5em]
\item Weil--Deligne 表示 $=$（$W_K$ 的连续表示 $\rho$）$+$（幂零算子 $N$，带 $q^{-\deg}$ 的
      共轭关系）。$N$ 必然幂零，且 $\varphi$-半单情形下 $N=0$ 恰对应「单」。
\item 任何 $\varphi$-半单的 $(\rho,N)$ 都是 $\bigoplus\mathrm{Sp}_m(r)$；$\mathrm{Sp}_m$ 是
      把 $N$ 的链长 $m$ 与 Tate 扭转 $\chi^{-i}$ 捆在一起的那个基本构件。
\item $l$-进表示与 Weil--Deligne 表示由 $\mathrm{WD}_{\zeta,\varphi}$ 等价；
      不可约 $l$-进表示 $\Leftrightarrow$ $N=0$ 且 $\rho'$ 不可约。
\item 惯性类型（Weil--Deligne 型 / Galois 型）只保留限制到 $I_K$ 的信息：
      前者忘掉 $N$，后者保留 Frobenius 的半单作用。第 4 节的惯性 Langlands 对应
      正是在这两类数据上做文章。
\end{itemize}
\end{keytakeaways}

% =============================================================================
%  §2 GLn 的表示（底本 p.10--22）
% =============================================================================
\section{$\GL_n$ 的表示}

% 2.1 容许表示（底本 p.10--16）
\subsection{$\GL_n$ 的容许表示}

\subsubsection{一般定义}

设 $\mathbf{G}$ 是 $K$ 上的连通约化群，记 $G=\mathbf{G}(K)$。设 $V$ 是 $\mathbb{C}$-向量空间
（下文中 $\mathbb{C}$ 可换成任意代数闭的特征 $0$ 域：事实上我们固定一个同构
$\iota\colon L=\overline{\mathbb{Q}}_l\xrightarrow{\ \sim\ }\mathbb{C}$，此后在笔记中不再区分这些系数域）。

\begin{definition}
设 $(\pi,V)$ 是 $G$ 的表示。
\begin{enumerate}[label=(\arabic*),leftmargin=2.2em]
\item $\pi$ 称为\emph{光滑的}，若每个 $v\in V$ 的稳定化子在 $G$ 中开。
\item $\pi$ 称为\emph{容许的}，若它光滑，且对每个开子群 $H\subset G$，$V^{H}$ 有限维。
\end{enumerate}
\end{definition}

记 $\Repinf_G$ 为光滑表示的范畴，$\Repadm_G$ 为容许表示的范畴。

\begin{remark}
注意 $\pi$ 光滑当且仅当 $V=\bigcup_K V^{K}$，其中 $K$ 遍历 $G$ 的所有紧开子群。
若进一步每个 $V^{K}$ 有限维，则 $\pi$ 容许。
\end{remark}

\begin{remark}
设 $K$ 是 $G$ 的任一紧开子群，$E(K)$ 是光滑不可约表示的同构类集合
（注意光滑性蕴含这些表示必有限维）。限制 $V|_K$ 显然半单：由光滑性，每个 $v\in V$
都落在某个（有限维的）$K$ 的不可约表示里。对每个 $\sigma=(\Sigma,\sigma)\in E(K)$，
记 $V_{\sigma}$ 为 $V$ 的 $\sigma$-等型分量，即 $V$ 中承载该类表示的最小 $K$-不变子空间；
于是 $V=\bigoplus_{\sigma\in E(K)}V_{\sigma}$，且每个 $V_{\sigma}|_K$ 是 $K$ 的表示 $\Sigma$
的直和。$(\pi,V)$ 容许，恰当每个 $V_{\sigma}$ 有限维；此时
\[
V|_K=\bigoplus V_{\sigma}=\bigoplus_{\sigma\in E(K)}\Sigma^{m_{\sigma}},
\]
其中 $m_{\sigma}\in\mathbb{Z}_{\ge 0}$ 是 $\sigma$ 在 $V|_K$ 中的重数。
\end{remark}

\begin{definition}
设 $(\pi,V)\in\Repinf_G$，$V^{*}$ 是 $V$ 的线性对偶（在向量空间范畴中），它按通常方式
带有 $G$ 的作用 $(g\psi)(v)=\psi(g^{-1}v)$。定义
\[
V^{\vee}=\widetilde{V}:=\{\psi\in V^{*}:\ \psi\ \text{的每个稳定化子在}\ G\ \text{中开}\}.
\]
$V^{\vee}=\widetilde{V}$ 称为表示 $V$ 的\emph{反变对偶}（contragradient），也记作
$\pi^{\vee}$ 或 $\widetilde{\pi}$；按定义它是光滑的。
\end{definition}

\begin{remark}
一般地，自然态射 $V\to V^{\vee\vee}$ 不是同构。从某种意义上说，容许性正是为了让这个
问题消失而定义的：
\begin{enumerate}[label=(\arabic*),leftmargin=2.2em]
\item $\pi$ 容许当且仅当 $\pi^{\vee}$ 容许；
\item $\pi$ 容许当且仅当 $V\xrightarrow{\ \sim\ }V^{\vee\vee}$ 是同构；
\item $\pi$ 不可约当且仅当 $\pi^{\vee}$ 不可约；
\item （Gelfand--Kazhdan）当 $G=\GL_n(K)$、$\pi\in\Repinf_G$ 不可约时，$\pi^{\vee}$ 同构于
      表示 $g\mapsto\pi({}^{t}g^{-1})$，这正是「反变对偶」这个名字的来历。
\end{enumerate}
\end{remark}

\subsubsection{Hecke 代数}

\begin{definition}
设 $G$ 如上，把 $G$ 上的 Haar 测度正规化，使 $\mu(K_0)=1$，其中 $K_0\subset G$ 是极大紧子群。
\begin{enumerate}[label=(\arabic*),leftmargin=2.2em]
\item \emph{Hecke 代数} $H(G)$ 是紧支集的局部常值函数 $f\colon G\to\mathbb{C}$ 构成的集合，
      在卷积下（取定 Haar 测度后）成为结合代数：
\[
(f_1*f_2)(h)=\int_G f_1(hg^{-1})f_2(g)\,dg=\int_G f_1(g)f_2(g^{-1}h)\,dg .
\]
\item 设 $K\subset G$ 是紧开子群，定义 $H(G/\!/K)$ 为 $H(G)$ 中在 $K$ 下左右不变的
      函数构成的子代数。
\item 更一般地，设 $Z\subset Z(G)$ 是中心的闭子空间，$\chi\colon Z\to\mathbb{C}$ 是 $Z$ 的
      特征标。有如下双陪集 Hecke 代数的变体：
\[
H_{\chi}(G):=\bigl\{f\colon G\to\mathbb{C}\ \text{局部常值、模}\ Z\ \text{紧支集}:\ 
f(zg)=\chi(z)^{-1}f(g)\ \ \forall z\in Z\bigr\},
\]
其乘法为 $(f_1*f_2)(h)=\int_{G/Z}f_1(x)f_2(x^{-1}h)\,dx$。
\end{enumerate}
\end{definition}

\begin{remark}
\leavevmode
\begin{enumerate}[label=(\arabic*),leftmargin=2.2em]
\item $H(G/\!/K)=\{\varphi\colon K\backslash G/K\to\mathbb{C}\ \text{有限支集}\}$。换言之，
      $H(G,K)$ 由那些满足 $\varphi(kgk')=\varphi(g)$（对一切 $k,k'\in K$）且支集为有限个
      双陪集 $KgK$ 之并的 $\varphi\in H(G)$ 构成。
\item $H(G/\!/K)$ 有双边单位元，即 $e_K:=\frac{1}{\vol(K)}\mathbf{1}_K$，$K$ 上正规化的
      特征函数。
\item 然而 $H(G)$ 不是含幺代数。$H(G)=\bigcup_K H(G/\!/K)$ 是各子代数之并，
      但对不同的 $K$ 单位元 $e_K$ 也不同。
\end{enumerate}
\end{remark}

\subsubsection{$G$ 在 $H(G)$ 上的作用}

$G$ 在 $H(G)$ 上有自然的左作用（相应地，右作用）\en{1}，现描述如下。
\begin{itemize}[leftmargin=1.6em]
\item $G$ 在 $H(G)$ 上的左作用记为 $\lambda_g$，由
\[
\lambda_gf(z)=f(g^{-1}z)\qquad(\forall z\in G)
\]
给出。显然对任意 $f\in H(G)$ 有 $\lambda_gf\in H(G)$，且这确实定义了作用。
\item 同样，右作用 $\rho_g$ 由 $\rho_gf(z)=f(zg)$（$\forall z\in G$）给出。
\end{itemize}

\begin{remark}
我们把 $H(G)$ 定义成 $C^{\infty}_c(G,\mathbb{C})$ 的部分原因，正是要使上面两个 $G$-作用
都成为 $G$ 的光滑表示。
\end{remark}

\subsubsection{Hecke 模}

设 $(\pi,V)$ 是光滑表示，$H(G)$ 通过
\[
f\star v:=\pi(f)(v)=\int_G f(g)\pi(g)v\,dg
\]
作用在 $V$ 上。由上面的注记，这仍是有限和。由于 $V=\bigcup_K V^{K}$，Hecke 作用使 $V$
成为非退化的 $H(G)$-模，即 $H(G)\cdot V=V$（例如 $\frac{1}{\mu(K)}\mathbf{1}_K\star v
=\pi(e_K)v=v$）。

\begin{lemma}
存在自然的函子
\[
\Repinf_G\longrightarrow\{\text{非退化}\ H(G)\text{-模}\},
\]
它诱导出范畴等价。
\end{lemma}

在继续之前，我们固定一些记号并记录几条简单观察。

\begin{enumerate}[label=(\arabic*),leftmargin=2.2em]
\item 对任意紧开子群 $K\subset G$，令 $e_{gK}:=\frac{1}{\mu(K)}\mathbf{1}_{gK}$ 为 $gK$ 上
      正规化的特征函数。注意当 $g=\mathrm{id}$ 时，它与 $e_K$ 的记号一致。
\item 直接验证可知：对 $f\in H(G)$，$e_K*f=f$ 当且仅当对一切 $k\in K$ 有 $f(kg)=f(g)$。
      特别地，对任意紧开子群 $K'\subset K$，有 $e_K=e_{K'}*e_K$。
\item 仍设 $K'\subset K$ 如上，$g\in G$ 任意。我们断言 $e_{gK}=e_{gK'}*e_K$。
      由定义，
\[
e_{gK'}*e_K(z)=\int_G e_{gK'}(x)\,e_K(x^{-1}z)\,dx,
\]
它非零当且仅当 $z\in xK$、其中 $x\in gK'$，即当且仅当 $z\in gK$。对这样的 $z$，
$e_{gK'}*e_K(z)$ 的值恰为
$\frac{1}{\mu(K)}\mathbf{1}_{K'}\cdot\mu(gK'\cap gK)=\frac{1}{\mu(K)}$。
此处关键在于第二个特征函数支撑在一个\emph{子群} $K$ 上（而非它的某个平移 $gK$ 上）。
\item 在上面描述的函子之下，$(\pi,V)$ 的子表示对应于 $H(G)$ 作用下的子模。
\end{enumerate}

\begin{proof}[引理 2.1 的证明]
我们已经构造了从 $\Repinf_G$ 到非退化 Hecke 模的函子；展开定义可以验证它全忠实。
现构造反方向的函子。设 $M$ 是非退化 Hecke 模。首先断言：对任意 $m\in M$，存在紧开子群
$K\subset G$ 使 $e_K\star m=m$。事实上，由非退化性可写 $m=\sum f_i\star m_i$（有限和，
$f_i\in H(G)$，$m_i\in M$）。取 $K$ 足够小，使 $K\subset G$ 在 $H(G)$ 上的左右作用都稳定
每个 $f_i$（见 2.1.3 小节），即 $f_i\in H(G/\!/K)$。于是
\[
e_K\star m=e_K\star\Bigl(\sum f_i\star m_i\Bigr)
=\sum\bigl((e_K*f_i)\star m_i\bigr)=m .
\]
取这样的紧开子群 $K$，定义
\[
g\cdot m=\pi(g)m:=e_{gK}\star m .
\]
首先这定义良好，即不依赖于 $K$ 的选取（只要它足够小使得 $e_K\star m=m$）。这很直接：
设 $K'\subset K$，则
\[
e_{gK'}\star m=e_{gK'}\star(e_K\star m)=(e_{gK'}*e_K)\star m=e_{gK}\star m,
\]
最后一步用的是上面的观察 (3)。

接着要证明这确实定义了作用，这一点其实颇为微妙（但光滑性是直接的）。
设 $g_1,g_2\in G$、$m\in M$，我们要选取 $K$ 与 $K'$ 使
\[
e_{g_1g_2K}\star m=e_{g_1K'}\star e_{g_2K}\star m
=(e_{g_1K'}*e_{g_2K})\star m .
\]
$K$ 的选取如前，使得 $e_K\star m=m$。现在令 $m':=g_2\cdot m$，断言可取 $K'$ 足够小，
使得 $e_{g_2K}\in H(G/\!/K')$ 且 $e_{K'}\star m'=m'$。为此注意 $e_{g_2K}\in H(G/\!/K')$
意即 $K'g_2K=g_2K$；先取 $K'=g_2Kg_2^{-1}\cap K$\en{2}，必要时再缩小以稳定 $m'$。
由此可直接验证
\[
e_{g_1K'}*e_{g_2K}=e_{g_1g_2K},
\]
从而确有 $e_{g_1g_2K}\star m=(e_{g_1K'}*e_{g_2K})\star m$；再把 $K$ 与 $K'$ 换成它们的交，
即得 $(g_1g_2)\cdot m=g_1\cdot(g_2\cdot m)$。
\end{proof}

\begin{remark}
同一构造还有另一种（更简洁的）做法：考虑 $H(G)$-模的自然映射
\[
H(G)\otimes_{H(G)}M\longrightarrow M,
\]
其中把 $H(G)$ 看作带右作用在其自身上的模。该映射其实是同构（尽管 $H(G)$ 不含单位元），
这来自非退化假设。于是通过同构 $H(G)\otimes M\cong M$、并让 $G$ 左作用在 $H(G)$ 上，
就得到 $M$ 上的 $G$-作用。注意这恰好对应
$g\cdot m=g\cdot(e_K\otimes m)=e_{gK}\otimes m$。
\end{remark}

\subsubsection{}

现在固定 $K$，则上面的函子限制为从 $\Repadm_G$ 到有限维 $H(G/\!/K)$-模范畴的函子。
在本书中，「不可约模」与「单模」作为同义词使用。

\begin{lemma}[Casselman]
固定紧开子群 $K\subset G$。
\begin{enumerate}[label=(\arabic*),leftmargin=2.2em]
\item 函子
\[
\bigl\{\pi\in\Rep^{f.l.}_{G}{}^{\mathrm{adm}}:\ 
\forall W\ \text{不可约子商},\ W^{K}\neq 0\bigr\}\xrightarrow{\ \sim\ }H(G/\!/K)\text{-}\mathrm{mod}^{f.d.},
\qquad V\longmapsto V^{K}
\]
是从「所有不可约子商都含 $K$-不动向量的有限长容许表示」范畴到有限维
$H(G/\!/K)$-模范畴的等价。
\item 若 $V\in\Repinf_G$ 不可约且 $V^{K}\neq 0$，则 $V^{K}$ 是不可约 $H(G/\!/K)$-模。
      此外，下述两类对象之间存在双射：
\begin{enumerate}[label=(\alph*),leftmargin=2.0em]
\item 满足 $V^{K}\neq 0$ 的不可约光滑表示 $(\pi,V)$ 的等价类；
\item 不可约 $H(G/\!/K)$-模的同构类。
\end{enumerate}
\end{enumerate}
\end{lemma}

\begin{remark}
容易证明：若 $V\in\Repinf_G$ 不可约，则 $V^{K}$ 或是 $0$、或是不可约 $H(G)$-模。
但反过来则是一个非平凡的陈述。
\end{remark}

作为推论，我们有

\begin{corollary}
若 $V\in\Repinf_G$ 不可约且 $H(G/\!/K)$ 交换，则或者 $V^{K}=0$，或者 $V^{K}$ 是一维的。
\end{corollary}

\subsubsection{球面 Hecke 代数}

在我们的设定里 $G=\GL_n(K)$。若 $C=\GL_n(\mathcal{O}_K)$ 是 $G$ 的极大紧开子群，
则 $H(G/\!/C)$ 交换；事实上
\[
H(G/\!/C)\xrightarrow{\ \sim\ }\mathbb{C}[t_1^{\pm 1},\dots,t_n^{\pm 1}]^{S_n}.
\]
更一般地，对 $G=\mathbf{G}(K)$（$\mathbf{G}$ 分裂约化，固定 $\mathbf{T}\subset\mathbf{B}\subset\mathbf{G}$）
与超特殊子群 $C\subset G$，经典的 Satake 同构说
\[
H(G/\!/C)\xrightarrow{\ \sim\ }(\mathcal{O}_{\widehat{T}})^{W},
\]
其中 $\widehat{T}=\Spec\mathbb{C}[\Lambda]$，$\Lambda$ 是余权格，$W$ 是 Weyl 群。
特别地，球面 Hecke 代数交换。

\begin{remark}
在 $\GL_n$ 情形，$H$ 的交换性也可以从 Gelfand 技巧看出。注意
\[
C\backslash G/C\cong G\backslash\bigl(G/C\times G/C\bigr),
\qquad [g]\longmapsto(1,[g]).
\]
而 $G/C$ 参数化 $K^{n}$ 中的格。也就是说可以等同
\[
C\backslash G/C\ \longleftrightarrow\ G\backslash(\mathcal{L}^{n}\times\mathcal{L}^{n})
=\{\text{格的相对位置}\},
\]
其中用 $\mathcal{L}^{n}$ 记 $K^{n}$ 中的格。从这个观点看，
\[
(\varphi*\psi)(\Lambda,\Lambda')=\sum_{\Lambda''}\varphi(\Lambda,\Lambda'')\psi(\Lambda'',\Lambda').
\]
注意此和有限，因为支集紧。Gelfand 的想法是：$(\Lambda,\Lambda')$ 的相对位置与
$((\Lambda')^{\vee},\Lambda^{\vee})$ 相同\en{3}，因此
\begin{multline*}
(\varphi*\psi)(\Lambda,\Lambda')
=\sum_{\Lambda''}\varphi(\Lambda,\Lambda'')\psi(\Lambda'',\Lambda')
=\sum_{\Lambda''}\varphi\bigl((\Lambda'')^{\vee},\Lambda^{\vee}\bigr)
\psi\bigl((\Lambda')^{\vee},(\Lambda'')^{\vee}\bigr)\\
=\sum_{\Lambda''}\varphi(\Lambda'',\Lambda^{\vee})\psi\bigl((\Lambda')^{\vee},\Lambda''\bigr)
=(\psi*\varphi)\bigl((\Lambda')^{\vee},\Lambda^{\vee}\bigr)=(\psi*\varphi)(\Lambda,\Lambda').
\end{multline*}
这个证明几乎是同义反复。
\end{remark}

\begin{corollarydefinition}
设 $(\pi,V)$ 是 $G$ 的不可约容许表示，$C=\GL_n(\mathcal{O}_K)$ 是 $G$ 的极大紧子群，
则 $V^{C}$ 或为 $0$、或为一维。若 $V^{C}\neq 0$（底本此处误作 $V^{\pi}\neq 0$，见译注），
则称 $V$ 是\emph{非分歧的}，即 $V^{C}$ 一维。
\end{corollarydefinition}

\begin{yizhu}[底本的笔误：$V^{\pi}$ 应为 $V^{C}$]
底本 2.1.6 末写作「$V$ is called unramified if $V^{\pi}\neq 0$, i.e., $V^{\pi}$ is
1-dimensional」。此处 $\pi$ 应是 $C$：全文用 $V^{C}$ 记 $C$-不动向量，
而 $V^{\pi}$ 没有意义（$\pi$ 是表示本身）。译文按 $V^{C}$ 处理，并在正文括注原文。
\end{yizhu}

\subsubsection{Schur 引理}

\begin{lemma}
设 $(\pi,V)\in\Repinf_G$ 是 $G$ 的光滑不可约表示，则 $V$ 的每个 $G$-自同态都是数乘。
\end{lemma}

与 $V$ 有限维的情形不同，Schur 引理在这里是非平凡的陈述。当
$\pi\in\Repadm_G$ 不可约且容许\en{4}时，可以借引理 2.2 回到有限维模的情形：
取 $C\subset G$ 足够小使 $V^{C}\neq 0$，则 $V^{C}$ 的 $H(G/\!/C)$-线性自同态只有数乘。

\begin{proof}[略证]
在我们的设定 $G=\mathbf{G}(K)$ 下，注意 $G$ 有可数基，且对每个紧开子群 $K\subset G$，
$H(G/\!/K)$ 在 $\mathbb{C}$ 上可数维，故 $H(G)$ 在 $\mathbb{C}$ 上可数维。
取 $V$ 中非零向量 $v$，由 $(\pi,V)$ 不可约，映射 $H(G)\to V$、$f\mapsto f\star v$ 是满射，
故 $V$ 可数维。再考虑环 $\End_G(V)$：仍由 $(\pi,V)$ 的不可约性，
映射 $\End_G(V)\to V$、$\mu\mapsto\mu(v)$ 是单射。于是 $\End_G(V)$ 是 $\mathbb{C}$ 上可数维的
除环；但不可数代数闭域 $F$ 上可数维的除环必为 $F$ 本身（证明：若 $D\neq F$，
取 $x\in D\setminus F$，则 $x$ 在 $F$ 上超越——因为 $F$ 代数闭；而元素族
$\{\frac{1}{x-\alpha}:\alpha\in F\}$ 不可数，且可验证其在 $F$ 上线性无关。证毕），
于是引理成立。
\end{proof}

\subsubsection{光滑不可约表示}

\begin{theorem}
设 $(\pi,V)$ 是 $G$ 的光滑不可约表示，则 $\pi$ 容许。
\end{theorem}

证明用 Jacquet 函子把 $\pi$ 嵌入到容许（超尖点）表示的抛物诱导中。
我们略去证明，请读者参阅 Jacquet 的处理或 [Bernstein--Zelevinsky]。

\begin{lemma}
$G=\GL_n(K)$ 的每个光滑不可约表示或是一维的、或是无限维的。
\end{lemma}

\begin{proof}
这来自 $G$ 的结构。设 $\pi\colon G\to\GL_n(\mathbb{C})$ 光滑，则
$\ker(\pi)\triangleleft G$ 是开子群，因而包含某个 $1+\varpi^{m}\mathfrak{M}_n(\mathcal{O}_K)
\cong 1+\varpi^{m}M_n(\mathcal{O}_K)$，从而包含 $1+\varpi^{m}\mathfrak{n}$，
即对角元为 $1$、其余元素落在 $\varpi^{m}\mathcal{O}_K$ 中的上三角矩阵全体。
由共轭可知幂幺上三角矩阵 $U\subset\ker\pi$。最后由正规性知 $\SL_n(K)\subset\ker\pi$。
这样我们就证明了引理，并得到更多：任何这样的有限维（从而一维）不可约表示 $\pi$
都通过行列式映射分解，即形如 $\chi\circ\det$，其中 $\chi$ 是 $K^{\times}$ 的某个拟特征标。
\end{proof}

\subsubsection{Harish-Chandra 特征标}\en{5}

回忆：给定光滑表示 $(\pi,V)$，Hecke 代数 $H(G)$ 通过
\[
\pi(f)(v)=\int_G f(g)\pi(g)v\,dg,\qquad f\in H(G)=C^{\infty}_c(G)
\]
作用在 $V$ 上。在「$G$ 的光滑表示 $\leftrightarrow$ 非退化 $H(G)$-模」的等价之下，
容许表示对应于那些「对一切 $f$，$\pi(f)$ 是有限秩算子」的 $H(G)$-模。
这使我们可以定义 $\pi(f)$ 的迹，从而定义分布特征标
\[
\Theta_{\pi}\colon H(G)\to\mathbb{C},\qquad f\longmapsto\Theta_{\pi}(f):=\Tr\pi(f).
\]

\begin{lemma}
若 $\{\pi_1,\dots,\pi_n\}$ 是 $G$ 的两两不同构的不可约容许表示，则
$\{\Theta_{\pi_1},\dots,\Theta_{\pi_n}\}$ 线性无关。
\end{lemma}

\begin{remark}
计算分布特征标 $\Theta_{\pi}$（$(\pi,V)$ 给定）有一个有用的办法如下。
取 $G$ 的一列可数基 $K_m$（例如 $K_i=1+\varpi^{i}K_0$，$K_0=\GL_2(\mathcal{O}_K)$\en{原注}），
令 $V_i:=V^{K_i}$，则 $V_0\subset V_1\subset\cdots\subset V=\bigcup_iV_i$。对每个 $i$ 与
每个 $g\in G$，考虑复合
\[
\pi(e_{K_i})\circ\pi(g)\circ\pi(e_{K_i})\colon V_i\longrightarrow V_i,
\]
记其迹为 $\xi_{\pi,i}(g)$。对 $f\in H(G)$，存在某个 $i$（依赖于 $f$）使 $f\in H(G/\!/K_i)$；
于是对每个 $m\ge i$ 有
\[
\Theta_{\pi}(f)=\int_G\xi_{\pi,m}(g)f(g)\,dg,
\]
且右端在 $m$ 足够大时不依赖于 $m$。
\end{remark}

我们陈述 Harish-Chandra 的一个深刻定理（事实上，当 $K$ 特征为 $0$ 时这是 Harish-Chandra
的定理，正特征时是 Lemaire 的定理）：$\Theta_{\pi}$ 可由 Harish-Chandra 特征标
$\chi_{\pi}$ 表示，后者在正则半单元 $G^{\mathrm{r.s.}}$ 上局部常值。先建立记号：
\begin{enumerate}[label=(\arabic*),leftmargin=2.2em]
\item 记 $G^{\mathrm{r.s.}}$ 为 $\GL_n(K)$ 中正则半单元的集合。（回忆 $g\in\GL_n(K)$
      正则半单，指其特征多项式 $P_g(X)\in k[X]$ 在 $K$ 上无重根。）
\item 记 $G^{\mathrm{r.e.}}\subset G^{\mathrm{r.s.}}$ 为正则椭圆元的子集。（回忆
      $g\in G^{\mathrm{r.s.}}$ 正则椭圆，指此外其特征多项式 $P_g$ 不可约。）
\end{enumerate}

\begin{theorem}[Harish-Chandra]
存在 $G$ 上的局部可积函数 $\chi_{\pi}$，它在正则半单集 $G^{\mathrm{r.s.}}$ 上局部常值、
在 $G^{\mathrm{r.s.}}$ 的补集上为 $0$，使得
\[
\Theta_{\pi}(f)=\int_G f(g)\chi_{\pi}(g)\,dg
\]
对一切 $f\in H(G)$ 成立。
\end{theorem}

\begin{remark}
$\chi_{\pi}$ 可以看作特征标 $\chi_{\pi}\colon\widetilde{G}\to\mathbb{C}$，其中
$\widetilde{G}=G^{\mathrm{r.s.}}/\sim$ 是 $G$ 中正则半单元的共轭类集合。
\end{remark}

\begin{remark}
我们记录一个在讨论 Jacquet--Langlands 对应时有用的定义。
\begin{enumerate}[label=(\arabic*),leftmargin=2.2em]
\item 光滑容许表示 $(\pi,V)$ 称为\emph{椭圆的}，若其 Harish-Chandra 特征标
      $\chi_{\pi}$ 在 $G^{\mathrm{r.e.}}$ 上不恒为 $0$。
\item 每个本质平方可积表示都是椭圆的。
\end{enumerate}
\end{remark}

% 2.2 超尖点表示（底本 p.16--20）
\subsection{超尖点表示}

\subsubsection{诱导与紧诱导}

我们在稍一般的情形下工作：设 $G$ 是局部紧（可以是完全不连通）拓扑群，
$H<G$ 是闭子群，考虑它们的光滑复表示。

\begin{definition}
定义两个函子 $\Ind,\ \ind\colon\Repinf_H\to\Repinf_G$：
\begin{enumerate}[label=(\arabic*),leftmargin=2.2em]
\item $\Ind^G_H(W):=\bigl\{f\colon G\to W\ \big|\ f(hx)=hf(x)\ \forall h\in H,\,x\in G;
\ \text{且}\ f\ \text{一致局部常值}\bigr\}$，
其中最后一个条件指：存在开子群 $C\subset G$（可依赖于 $x$）使 $f(xC)=f(x)$ 对一切 $x$ 成立。
$G$ 通过右平移作用在 $\Ind(W)$ 上。
\item $\ind^G_H\subset\Ind^G_H$ 是子函子，
\[
\ind(W)=\{f\in\Ind(W)\mid f\ \text{模}\ H\ \text{紧支集}\}.
\]
\end{enumerate}
\end{definition}

\begin{remark}
\leavevmode
\begin{enumerate}[label=(\arabic*),leftmargin=2.2em]
\item 若 $G/H$ 紧，则 $\ind=\Ind$。
\item 设 $G/H$ 紧，则 $\ind=\Ind$ 把 $H$ 的容许表示送到 $G$ 的容许表示。
\item $\Ind$ 正合（$\Res$ 限制函子也正合）；若进一步 $H<G$ 是开子群，则 $\ind$ 也正合。
\end{enumerate}
\end{remark}

考虑限制函子 $\Res\colon\Repinf_G\to\Repinf_H$。

\begin{lemma}
\leavevmode
\begin{enumerate}[label=(\arabic*),leftmargin=2.2em]
\item （Frobenius 互反）\en{6} 有伴随对 $\Res\dashv\Ind$。
\item 进一步设 $H<G$ 是开子群（于是 $G/H$ 离散），则 $\ind\dashv\Res$。
\end{enumerate}
\end{lemma}

\begin{remark}
\leavevmode
\begin{enumerate}[label=(\arabic*),leftmargin=2.2em]
\item 函子 $\ind$、$\Res$、$\Ind$ 在某种意义上对应于 $f_!$、$f^{*}$、$f_{*}$，
      这可以严格化。[例如参见 Ngo 的笔记] 记住上面的伴随对时，这个类比很有用。
\item 设 $H<G$ 是开子群，则 $\ind$ 可以解释为 $\mathrm{coInd}$，即
      $\ind(W)=W\otimes_{\mathbb{C}[H]}\mathbb{C}[G]$。
\end{enumerate}
\end{remark}

现在考虑商 $\psi\colon G\to M=G/N$，其中 $N$ 是某个（闭）正规子群。我们仍有「$\Res$」函子：
\[
\Res=\psi^{*}\colon\Repinf_M\longrightarrow\Repinf_G,
\]
由与 $\psi$ 复合给出（有时称为\emph{膨胀函子}）。$\Res$ 的左伴随是取 $N$-余不变量子空间，
$\Res$ 的右伴随是取 $N$-不变量。

我认为下面这个约定与一些已有的约定不同：

\begin{definition}
设 $\psi\colon G\to M=G/N$ 如上，定义
\[
\psi_{!}\colon\Repinf_G\to\Repinf_M,\qquad \psi_{!}V=V/[N-1]V,
\]
即 $V$ 的 $N$-余不变量空间。
\end{definition}

于是如所期望，在此设定下我们有伴随对 $\psi_{!}\dashv\psi^{*}$。

\subsubsection{抛物诱导与 Jacquet 函子 I}

设 $P<G=\GL_n(K)$ 是抛物子群\en{7}，与 $P$ 相伴有幂幺根 $U_P$ 与 Levi 商 $L_P$，
它们适合正合列
\[
1\longrightarrow U_P\longrightarrow P\longrightarrow L_P\longrightarrow 1 .
\]
我们有 $G=\GL_n(K)$，故抛物子群都是标准抛物子群的共轭：
\[
P_{\underline n}=
\begin{pmatrix}
\GL_{n_1} & * & *\\
 & \GL_{n_2} & *\\
 & & \ddots
\end{pmatrix},
\qquad \underline n=(n_1,\dots,n_r),
\]
其中 $\underline n$ 是 $n$ 的一个有序分拆。在 $P_{\underline n}$ 内有
\[
U_{\underline n}=
\begin{pmatrix}
1_{n_1} & * & *\\
 & 1_{n_2} & *\\
 & & \ddots
\end{pmatrix},
\qquad
L_{\underline n}=
\begin{pmatrix}
\GL_{n_1} & & \\
 & \ddots & \\
 & & \GL_{n_r}
\end{pmatrix}.
\]
现在对任一标准抛物子群 $P$（$P<G$，$\phi\colon P\twoheadrightarrow L=P/U$），
由上面的讨论我们有函子
\[
\ind=\Ind\colon\Repinf_P\to\Repinf_G,\qquad
\phi_{!}\colon\Repinf_P\to\Repinf_L,
\]
以及伴随对
\[
\Res^G_P\dashv\ind^G_P=\Ind^G_P,\qquad \phi_{!}\dashv\phi^{*}.
\]
现在定义（朴素 / 未正规化的）抛物诱导 $\widetilde I^G_P$ 与 $\widetilde J^G_P$ 如下：

\begin{definition}
\leavevmode
\begin{enumerate}[label=(\arabic*),leftmargin=2.2em]
\item $\widetilde I^G_P\colon\Repinf_L\to\Repinf_G$ 定义为 $\widetilde I^G_P:=\Ind^G_P\circ\phi^{*}$；
\item $\widetilde J^G_P\colon\Repinf_G\to\Repinf_L$ 定义为
      $\widetilde J^G_P=\phi_{!}\circ\Res^G_P$。
\end{enumerate}
\end{definition}

\begin{remark}
我们有以下性质：
\begin{enumerate}[label=(\arabic*),leftmargin=2.2em]
\item $\widetilde J\dashv\widetilde I$；
\item 两个函子都正合；
\item 两个函子都把有限长表示送到有限长表示；
\item 两个函子都把容许表示送到容许表示；
\item 上述性质在用特征标扭转之后依然成立。
\end{enumerate}
其中 (1) 是 $\Res\dashv\Ind$ 与 $\phi_{!}\dashv\phi^{*}$ 的形式推论：
\[
\Hom_L(\widetilde J\sigma_G,\sigma_L)\cong
\Hom_P\bigl(\Res^G_P(\sigma_G),\phi^{*}(\sigma_L)\bigr)\cong
\Hom_G(\sigma_G,\widetilde I\sigma_L).
\]
\end{remark}

\subsubsection{抛物诱导与 Jacquet 函子 II}

朴素的 Jacquet 函子与抛物诱导并不把酉表示送到酉表示（因为抛物子群一般不是幺模的）。
设 $\delta$ 是 $P$ 上的模特征标，即由
\[
(L_x)^{*}\mu^{r}_P=\delta(x)\mu^{r}_P
\]
给出，其中 $L_x$ 是 $x$ 的左平移，$\mu^{r}_P$ 是 $P$ 上（任一）右不变 Haar 测度。

\begin{example}
考虑 $B\subset\GL_2(K)$ 是上三角矩阵构成的标准 Borel 子群
$\bigl\{\begin{psmallmatrix}x&y\\&z\end{psmallmatrix}\bigr\}\subset\GL_2$，
则其左不变测度 $\mu^{l}$ 与右不变测度 $\mu^{r}$ 分别为
\[
d\mu^{l}=\frac{1}{|x|^{2}\cdot|z|}\,dx\,dy\,dz,
\qquad
d\mu^{r}=\frac{1}{|x|\cdot|z|^{2}}\,dx\,dy\,dz .
\]
按定义 $d\mu^{r}(\gamma g)=\delta(\gamma)d\mu^{r}(g)$，于是容易算出
\begin{multline*}
d\mu^{r}\Bigl(\begin{pmatrix}a&b\\&d\end{pmatrix}
\begin{pmatrix}x&y\\&z\end{pmatrix}\Bigr)
=\frac{d(ax)\wedge d(ay+bz)\wedge d(dz)}{|ax|\cdot|dz|^{2}}\\
=\frac{|a|}{|d|}\cdot d\mu^{r}\Bigl(\begin{pmatrix}x&y\\&z\end{pmatrix}\Bigr),
\end{multline*}
故 $\delta\bigl(\begin{psmallmatrix}a&b\\&d\end{psmallmatrix}\bigr)=\bigl|a\cdot d^{-1}\bigr|$。

更一般地，对 $L_{\underline n}\subset P_{\underline n}$（把 Levi 看作子群），
有 $P_{\underline n}$ 上的模特征标 $\delta_{\underline n}$，而 $\delta_{\underline n}|_{L_{\underline n}}$
可由
\[
\delta_{\underline n}(x)=\bigl|\det\bigl(\ad_{U_{\underline n}}(x)\bigr)\bigr|,
\qquad \forall x\in L_{\underline n}
\]
描述。这里 $\ad_{U_{\underline n}}$ 记 $L_{\underline n}$ 通过共轭作用在 $U_{\underline n}$ 上所得的表示。
\end{example}

这段讨论导向正规化 Jacquet 函子 $J$ 与抛物诱导 $I$ 的定义。

\begin{definition}
\[
J^G_P:=\delta^{-\frac12}\cdot\widetilde J^G_P,
\qquad
I^G_P:=\widetilde I^G_P\cdot\delta^{\frac12}.
\]
\end{definition}

\begin{remark}
按定义，对 $L_{\underline n}=\GL_{\underline n}$ 的任一光滑表示 $(\pi_L,W)$，有
\begin{multline*}
I^G_{P_{\underline n}}(W)=
\Bigl\{f\colon G\to W\ \Big|\ f\ \text{光滑},\ 
f(uhg)=\delta_{\underline n}^{1/2}(h)\,\pi_L(h)f(g)\\
\forall u\in U_{\underline n},\ h\in\GL_{\underline n},\ g\in G\Bigr\}.
\end{multline*}
\end{remark}

\subsubsection{超尖点表示}

\begin{definition}
设 $(\pi,V)\in\Repadm_G$ 是不可约容许表示。若对任一真抛物子群 $P<G$ 都有
$J^G_P(\pi)=0$，则称 $(\pi,V)$ 是\emph{超尖点的}。
\end{definition}

超尖点表示构成容许表示的基本构件。

\begin{lemma}
设 $\pi\in\Repinf_G$ 不可约（从而容许），则或者 $\pi$ 超尖点，或者 $\pi$ 是
$I^G_P(\sigma)$ 的子表示，其中 $P<G$ 是（真）抛物子群，$\sigma$ 是 $L$ 的不可约超尖点表示。
\end{lemma}

\begin{proof}
设 $\pi$ 不超尖点，则存在极小抛物子群 $P<G$ 使 $J^G_P(\pi)\neq 0$；$J^G_P\pi$ 有限长，
取其不可约商 $\sigma$。由伴随性
\[
0\neq\Hom_L(J^G_P\pi,\sigma)=\Hom_G(\pi,I^G_P\sigma),
\]
即 $\pi$ 同构于 $I^G_P\sigma$ 的某个子表示。按构造可以证明 $\sigma$ 超尖点。
\end{proof}

下述超尖点表示的描述（我认为属于 Harish-Chandra）很有用。令
$G^{0}=\{g\in\GL_n(K):\det g\in\mathcal{O}_K^{\times}\}$，$Z=Z(G)$ 为 $G$ 的中心，
则 $ZG^{0}$ 在 $G$ 中有限指数（这对所有约化 $p$-进群都成立，对 $\GL_n$ 指数为 $n$）。
其想法是：通常可以先研究 $G$ 的表示在 $G^{0}$ 上的限制，再单独处理中心。
本书后文会给出这一想法的应用。

\begin{lemma}
设 $(\pi,V)$ 是 $G$ 的不可约容许表示，则下列等价：\en{8}
\begin{enumerate}[label=(\arabic*),leftmargin=2.2em]
\item $\pi$ 超尖点；
\item 每个矩阵系数 $C_{v,\lambda}\colon g\mapsto\lambda(\pi(g)v)$（$G\to\mathbb{C}$）
      在 $G^{0}$ 上紧支集；
\item 每个矩阵系数模中心 $Z\subset G$ 紧支集。
\end{enumerate}
此外，若再设 $\pi^{\vee}$ 也不可约，则 (1)--(3) 还等价于
\begin{enumerate}[label=(\arabic*),leftmargin=2.2em]\setcounter{enumi}{3}
\item 存在一个矩阵系数模中心紧支集。
\end{enumerate}
\end{lemma}

\begin{proof}[略证]
底本此处只写「Add sketch in second draft」，未给证明。
\end{proof}

下面这个引理对我们有用：

\begin{lemma}
设 $K_0$ 是 $\GL_n(K)$ 的紧开子群。则 $G^{0}$ 的不可约超尖点表示中，
含有 $K_0$-不动向量的同构类只有有限多个。
\end{lemma}

\begin{proof}
这是 Bernstein 哈佛讲义第 3 节命题 21 推论的改写，见 [1]。
\end{proof}

\begin{remark}
下述从 $L_{\underline n}$ 到 $G$ 的抛物诱导的描述，将是 Bernstein--Zelevinsky 分类的特例。
\begin{itemize}[leftmargin=1.6em]
\item 设 $(\pi_L,W)$ 是 $L_{\underline n}=\GL_{n_1}\times\cdots\times\GL_{n_r}$ 的容许表示，
      由
\[
\pi_{\underline n}=\pi_1\otimes\cdots\otimes\pi_r,\qquad
W=V_1\otimes\cdots\otimes V_r
\]
给出，其中每个 $(\pi_i,V_i)$ 是 $\GL_{n_i}$ 的不可约超尖点表示。则表示
\[
\pi_1\times\cdots\times\pi_r:=I^G_{P_{\underline n}}(\pi_L)
\]
是 $G$ 的容许表示（如上所述）。它不可约当且仅当对任意 $i\neq j$ 的 $\pi_i,\pi_j$，
$\pi_i\neq|\det|\cdot\pi_j$。
\item 在 $\GL_2$ 情形，这说的正是：$\chi_1\times\chi_2$ 不可约当且仅当
      $\chi_1/\chi_2\neq|\cdot|^{\pm 1}$——这些就是 $\GL_2$ 的主序列表示。
\end{itemize}
\end{remark}

% 2.3 局部 Langlands 对应的陈述（底本 p.21--22）
\subsection{局部 Langlands 对应的陈述}

到目前为止，我们还没有引入 Galois 表示与容许表示的 $L$-因子和 $\epsilon$-因子，
也没有引入（本质）平方可积、（本质）缓增与一般（generic）表示的概念。
不过我认为此刻已经可以陈述精确的局部 Langlands 对应（LLC）了。

\subsubsection{LLC 的陈述}

\begin{theorem}
存在唯一的一族集合间双射
\[
\mathrm{rec}_n\colon
\left\{\begin{array}{c}
\GL_n(K)\ \text{的不可约容许表示}\\
\text{（的等价类）}
\end{array}\right\}
\xrightarrow{\ \sim\ }
\left\{\begin{array}{c}
W_K\ \text{的}\ n\ \text{维 Frobenius 半单}\\
\text{Weil--Deligne 表示（的等价类）}
\end{array}\right\},
\]
满足下列诸条性质：
\begin{enumerate}[label=(\arabic*),leftmargin=2.2em]
\item 对 $K^{\times}$ 的每个不可约拟特征标 $\chi$，
\[
\mathrm{rec}_1(\chi)=\chi\circ\Art_K^{-1},
\]
其中 $\Art_K\colon K^{\times}\xrightarrow{\ \sim\ }W^{\mathrm{ab}}_K$。
\item 对 $\GL_{n_1}$、$\GL_{n_2}$ 的不可约容许表示 $\pi_1,\pi_2$，
\[
L(\pi_1\times\pi_2,s)=L\bigl(\mathrm{rec}_{n_1}(\pi_1)\otimes\mathrm{rec}_{n_2}(\pi_2),s\bigr),
\qquad
\epsilon(\pi_1\times\pi_2,s)=\epsilon\bigl(\mathrm{rec}_{n_1}(\pi_1)\otimes
\mathrm{rec}_{n_2}(\pi_2),s\bigr),
\]
其中 $L$-因子与 $\epsilon$-因子在第 3.6 小节中定义。
\item 对 $K^{\times}$ 的每个拟特征标 $\chi$ 与每个不可约 $\pi\in\Repadm_{\GL_n(K)}$，有
\[
\mathrm{rec}_n(\pi\chi)=\mathrm{rec}_n(\pi)\otimes\mathrm{rec}_1(\chi).
\]
\item 对每个中心特征标为 $\omega_{\pi}$ 的不可约 $\pi\in\Repadm_{\GL_n(K)}$，
\[
\det\circ\mathrm{rec}_n(\pi)=\mathrm{rec}_1(\omega_{\pi}).
\]
\item 对每个不可约 $\pi\in\Repadm_{\GL_n(K)}$，有
\[
\mathrm{rec}_n(\pi^{\vee})=\mathrm{rec}_n(\pi)^{\vee}.
\]
\end{enumerate}
\end{theorem}

\begin{remark}
局部 Langlands 对应 $\mathrm{rec}_n$ 保持两个重要的不变量：
\begin{itemize}[leftmargin=1.6em]
\item \emph{导子}：$c(\pi)=c(\mathrm{rec}_n(\pi))$，这来自与 $\epsilon$-因子的相容性。
\item \emph{深度}：$d(\pi)=d(\mathrm{rec}_n(\pi))$，由 Moy--Prasad 引入。
\end{itemize}
\end{remark}

\begin{remark}
由 Henniart 的工作，只需证明存在一族满足上述 (1)--(5) 的映射 $\mathrm{rec}_n$
（作用在超尖点表示上）。Henniart 的定理断言：一旦 $\mathrm{rec}_n$ 存在，
它必然是一族双射且唯一。Henniart 的策略大致如下：
\begin{enumerate}[label=(\roman*),leftmargin=2.2em]
\item 由 Bernstein--Zelevinsky（将在第 3 节讨论），可以归结为证明存在唯一一族双射
\[
\mathrm{rec}_n\colon\{\GL_n\ \text{超尖点表示的同构类}\}\longrightarrow
\{\text{$W_K$ 的}\ n\ \text{维不可约 WD 表示的同构类}\};
\]
\item 由 $\epsilon$-因子的性质（并同时考虑所有 $n$ 的 $\mathrm{rec}_n$），把上面的陈述
      归结为双射族 $\{\mathrm{rec}_n\}$ 的存在性；
\item 由 Henniart 的数值局部 Langlands 对应，映射 $\mathrm{rec}_n$ 的满性可由单射性推出；
\item 用 $L$-因子的性质证明这样的族 $\{\mathrm{rec}_n\}$ 必然单射。
\end{enumerate}
于是只需证明：在 $\GL_n$ 的不可约超尖点表示上存在满足定理中 (1)--(5) 的映射族
$\{\mathrm{rec}_n\}$。这正是 Harris--Taylor 证明、其后又由 Henniart 在 $K/\mathbb{Q}_p$
为局部数域时证明的陈述\en{译注}。看来 Scholze 的方法绕过了数值局部 Langlands 对应。
\end{remark}

\begin{yizhu}[底本的笔误：local number field]
底本此句作 “when $K/\mathbb{Q}_p$ is a local number field”，字面是「局部数域」，
但应与前一空格的 $\mathbb{Q}_p$ 一致，读作「$p$-进域」（$K/\mathbb{Q}_p$ 有限）。
译文保留「局部数域」并在方括号内标出正读。另：Harris--Taylor 与 Henniart 的原始结论
是对 $p$-进域上的 $\GL_n$ 的 LLC，这正是本文全篇的底域设定。
\end{yizhu}

\subsubsection{更精细的对应}

注意上面定理中包含的性质是为保证对应的唯一性所必需的（因而属于定理的一部分）。
下面几条可以看作额外的性质（它们对 $\mathrm{rec}_n$ 给出更精细的信息）：

\begin{theorem}
在上述局部 Langlands 对应（即定理中的 $\mathrm{rec}_n$）之下，有
\begin{enumerate}[label=(\arabic*),leftmargin=2.2em]
\item 不可约表示 $\pi\in\Repadm_G$ 超尖点，当且仅当 $\mathrm{rec}_n(\pi)$ 不可约；
\item $\pi$ 本质平方可积，当且仅当 $\mathrm{rec}_n(\pi)$ 不可分解；
\item $\pi$ 一般（generic），当且仅当 $L(\ad\circ\mathrm{rec}_n(\pi),s)$ 在 $s=1$ 处无极点。
\end{enumerate}
\end{theorem}

\subsubsection{惯性对应}

我们将在下一节陈述（$\GL_n(\mathcal{O}_K)$-）类型论的特例。虽然不能指望
$\GL_n(\mathcal{O}_K)$-表示与惯性表示之间存在对应，但两者之间有密切的联系。

\begin{proposition}
\leavevmode
\begin{enumerate}[label=(\arabic*),leftmargin=2.2em]
\item 设 $\pi$ 与 $\pi'$ 是 $G$ 的两个一般表示。若
\[
r_n(\pi)|_{I_K}\cong r_n(\pi')|_{I_K},
\]
其中 $r_n$ 记与 Weil--Deligne 表示 $\mathrm{rec}_n$ 相伴的 $W_K$ 的 $l$-进表示，则
\[
\pi|_{\GL_n(\mathcal{O}_K)}\cong\pi'|_{\GL_n(\mathcal{O}_K)}.
\]
\item （Paskunas）设 $\tau$ 是 $I_K$ 的 Weil--Deligne 惯性类型，且可延拓为 $W_K$ 的
      $n$ 维不可约 Frobenius 半单表示，则存在唯一（同构意义下）的光滑不可约表示
      $\sigma=\sigma_{\tau}$（$\GL_n(\mathcal{O}_K)$ 的），使得对 $\GL_n(K)$ 的任一光滑不可约
      无限维表示 $\pi$，
\[
\pi|_K\ \text{包含}\ \sigma_{\tau}
\quad\Longleftrightarrow\quad
\mathrm{rec}_n(\pi)|_{I_K}\cong\tau .
\]
\end{enumerate}
\end{proposition}

\begin{keytakeaways}
\begin{itemize}[leftmargin=1.5em]
\item LLC 就是一族双射 $\mathrm{rec}_n$：$\GL_n$ 的不可约容许表示 $\leftrightarrow$
      $W_K$ 的 $n$ 维 Frobenius 半单 Weil--Deligne 表示，由 (1)--(5) 五条性质唯一确定。
\item 唯一性由 Henniart 的论证保证：只要在超尖点表示上造出满足 (1)--(5) 的 $\mathrm{rec}_n$，
      其余由 Bernstein--Zelevinsky 分类自动拼装，且自动是双射。
\item 更精细的对应：超尖点 $\leftrightarrow$ 不可约；本质平方可积 $\leftrightarrow$ 不可分解；
      一般 $\leftrightarrow$ $L(\ad\circ\mathrm{rec},s)$ 在 $s=1$ 无极点。
\item 惯性对应说明 $\GL_n(\mathcal{O}_K)$ 上的表示（类型）与 $I_K$ 上的惯性类型之间有
      一一的联系，这是第 4 节的主题。
\end{itemize}
\end{keytakeaways}


% 3. Bernstein-Zelevinsky（及其他）分类（底本 p.23--31）
\section{Bernstein--Zelevinsky（及其他）分类}

\subsection{$\GL_n$ 容许表示的分类}

\subsubsection{段 $\Delta(\pi,m)$}

设 $(\pi,V)\in\Repadm_G$，对任意 $s\in\mathbb{C}$，记 $\pi(s)$ 为 $\pi$ 用特征标
$|\det|^{s}$ 所作的扭转，即 $\pi(s)(g)=|\det(g)|^{s}\pi(g)$。
注意若 $\pi$ 不可约超尖点，则 $\pi(s)$ 也是。

我们在 $\GL_n(K)$ 的不可约超尖点表示的等构类集合上给一个偏序：
$\pi\le\pi'$ 当且仅当存在整数 $s\ge 0$ 使 $\pi'=\pi(s)$。记这个偏序集为 $\Asc_n$。
$\Asc_n$ 中的有限区间称为\emph{Zelevinsky 段}。更精确地说：

\begin{definition}
从 $\pi$ 出发、长度为 $m$、次数为 $mn$ 的\emph{Zelevinsky 段}\en{9}是
$\GL_n(K)$ 的 $m$ 个不可约超尖点表示构成的有序有限族：
\[
\Delta(\pi,m)=[\pi,\pi(1),\dots,\pi(m-1)].
\]
\end{definition}

\begin{definition}
设 $\Delta_1=(\pi_1,m_1)$ 与 $\Delta_2=(\pi_2,m_2)$ 分别是 $\Asc_{n_1}$ 与 $\Asc_{n_2}$
中的两个 Zelevinsky 段。
\begin{enumerate}[label=(\arabic*),leftmargin=2.2em]
\item 若 $\Delta_1\not\subseteq\Delta_2$、$\Delta_2\not\subseteq\Delta_1$，且
      $\Delta_1\cup\Delta_2$ 是 Zelevinsky 段，则称 $\Delta_1$ 与 $\Delta_2$ \emph{相链}
      （linked）。特别地，这只有在 $n_1=n_2$ 时才可能发生。
\item 若 (a) $\Delta_1$ 与 $\Delta_2$ 相链，且 (b) $\pi_1<\pi_2$，则称
      $\Delta_1$ \emph{先于} $\Delta_2$（precedes）。
\end{enumerate}
\end{definition}

\begin{remark}
特别地，若 $n_1\neq n_2$，则 $\Delta_1$ 与 $\Delta_2$ 不相链。
\end{remark}

\begin{definition}
设 $\Delta=\Delta(\pi,m)$ 是（Zelevinsky）段，长度为 $m$、次数为 $mn$，
定义 $\GL_{mn}$ 的容许表示
\[
\pi(\Delta):=\pi_1\times\cdots\times\pi_m,
\]
其中 $\pi_1\times\cdots\times\pi_m=I^{\GL_{mn}}_{P_{(n,n,\dots,n)}}(\otimes\pi_i)$，
如 2.2.4 末的注所述。
\end{definition}

\subsubsection{$(*)$ 条件}

设 $\Delta_1,\dots,\Delta_r$ 是 $r$ 个长度分别为 $m_i$、位于 $\Asc_{n_i}$ 中的
Zelevinsky 段构成的有序族。引入下述 $(*)$ 条件：
\[
(*) \qquad \text{对每个 } i<j,\ \Delta_i\ \text{不先于}\ \Delta_j .
\]
$\Delta_1,\dots,\Delta_r$ 的一个 \emph{$(*)$-置换}是指 $\Delta_i$ 上的一个置换 $\sigma$，
使得 $\Delta_{\sigma(1)},\dots,\Delta_{\sigma(r)}$ 仍满足 $(*)$ 条件。

注意：
\begin{enumerate}[label=(\arabic*),leftmargin=2.2em]
\item 按定义，当 $n_i\neq n_j$ 时该条件为空；
\item $(*)$ 条件允许在 $i<j$ 时 $\Delta_j$ 先于 $\Delta_i$。
\end{enumerate}

\begin{example}
设 $\Delta=[\pi,\pi(1)]$，$\Delta'=[\pi(2),\pi(3)]$，$\Delta''=[\pi(5),\pi(6)]$，
则 $(\Delta,\Delta')$ 不满足 $(*)$，而 $(\Delta',\Delta)$、$(\Delta,\Delta'')$，
以及单个段 $\Delta\cup\Delta'=[\pi,\dots,\pi(3)]$ 都满足 $(*)$。
\end{example}

\subsubsection{分类定理的陈述}

\begin{theorem}[Bernstein--Zelevinsky]
\leavevmode
\begin{enumerate}[label=(\arabic*),leftmargin=2.2em]
\item 对每个长度为 $m$ 的 Zelevinsky 段 $\Delta$，表示 $\pi(\Delta)$ 的合成列长度为 $2^{m-1}$。
\item 设 $\Delta$ 是有限长的 Zelevinsky 段，则 $\pi(\Delta)$ 有唯一的不可约子表示 $Z(\Delta)$
      与唯一的不可约商表示 $Q(\Delta)$。
\item 设 $\Delta_1,\dots,\Delta_r$ 是满足 $(*)$ 条件的 Zelevinsky 段，则表示
\[
Q(\Delta_1)\times\cdots\times Q(\Delta_r)\in\Repadm_{\GL_{\sum n_im_i}}
\]
有唯一的不可约商 $Q(\Delta_1,\dots,\Delta_r)$。此外，在 $(*)$-置换之下，
$Q(\Delta_1,\dots,\Delta_r)$ 与 $\Delta_i$ 的次序无关。
\item $G=\GL_n(K)$ 的任一不可约表示 $(\pi,V)$ 都同构于如下形式：
\[
\pi\cong Q(\Delta_1,\dots,\Delta_r),
\qquad
\Delta_i=(\pi_i,m_i),\quad n=\sum n_im_i,
\]
其中段族 $\Delta_1,\dots,\Delta_r$ 在 $(*)$-共轭之下唯一。
\item 仍如 (2)，设 $\Delta_1,\dots,\Delta_r$ 是满足 $(*)$ 条件的段，则
\[
Q(\Delta_1)\times\cdots\times Q(\Delta_r)
\]
不可约（从而等于 $Q(\Delta_1,\dots,\Delta_r)$）当且仅当任两段 $\Delta_i$ 与 $\Delta_j$ 都不相链。
作为一个特例，我们重新得到 2.2.4 中的注。
\item 对 $\Delta=\Delta(\pi,m)=[\pi,\dots,\pi(m-1)]$，定义
\[
\Delta^{\vee}=[\pi(m-1)^{\vee},\dots,\pi^{\vee}]=\Delta\bigl(\pi^{\vee}(1-m),m\bigr),
\]
则
\[
Q(\Delta_1^{\vee},\dots,\Delta_r^{\vee})=Q(\Delta_1,\dots,\Delta_r)^{\vee}.
\]
\item (3)--(6) 的类似断言对 $Z(\Delta_1)\times\cdots\times Z(\Delta_r)$ 成立：
      它有唯一的不可约子表示 $Z(\Delta_1,\dots,\Delta_r)$。
\end{enumerate}
\end{theorem}

\begin{remark}
任给 Zelevinsky 段的有序族 $\Delta_1,\dots,\Delta_r$，总存在它的某个置换满足 $(*)$ 条件，
于是 $\GL_n(K)$ 的容许表示的等构类与无序族 $\Delta_1,\dots,\Delta_r$（$n=\sum m_in_i$）
一一对应。
\end{remark}

\begin{remark}
可以给各 $\GL_n(K)$ 的容许表示的 Grothendieck 群之直和 $R_K=\bigoplus_{n\ge 0}R_n(K)$
一个乘法结构
\[
([\pi_1],[\pi_2])\longmapsto[\pi_1\times\pi_2],
\]
使 $R_K$ 成为分次交换环。而映射 $Q(\Delta_1,\dots,\Delta_r)\mapsto
Z(\Delta_1,\dots,\Delta_r)$ 定义了 $R_K$ 上的一个对合。
\end{remark}

\subsubsection{超尖点支撑}

\begin{definition}
设 $\pi\cong Q(\Delta_1,\dots,\Delta_r)$ 如 3.1 所述，定义 $\pi$ 的\emph{超尖点支撑}为
下述（无序的）超尖点表示族：
\[
\operatorname{Supp}(\pi)=\bigl\{\pi_i(j)\bigr\}_{i=1,\dots,r;\ j=1,\dots,m_i-1}.
\]
\end{definition}

\begin{remark}[Zelevinsky]
$\pi$ 的超尖点支撑是唯一的一族无序超尖点表示 $\sigma_1,\dots,\sigma_s$，使得
$Q(\Delta_1,\dots,\Delta_r)$ 与 $Z(\Delta_1,\dots,\Delta_r)$ 都是
$\sigma_1\times\cdots\times\sigma_s$ 的子商。
\end{remark}

\subsubsection{Steinberg 表示}

考虑（拟）特征标 $|\cdot|^{\frac{1-n}{2}}\colon K^{\times}\to\mathbb{C}^{\times}$
以及 Zelevinsky 段
\[
\Delta=\Delta\Bigl(|\cdot|^{\frac{1-n}{2}},n\Bigr)
=\Bigl[|\cdot|^{\frac{1-n}{2}},|\cdot|^{\frac{3-n}{2}},\dots,|\cdot|^{\frac{n-1}{2}}\Bigr].
\]
直接计算标准 Borel 子群 $B$ 上的模特征标 $\delta_B$ 在 $T=L_B=\GL_{(1,\dots,1)}(K)$
上的限制。设 $x=\mathrm{Diag}(x_1,\dots,x_n)$ 是对角元为 $x_1,\dots,x_n$ 的对角矩阵，则
\[
\delta(x)=|\det(\ad_U(x))|
=\Bigl|x_1^{n-1}\cdot x_2^{n-3}\cdots x_n^{1-n}\Bigr|.
\]
换言之，对角环面 $T$ 上的表示 $\pi_T$（此时是一个特征标）
$|\cdot|^{\frac{1-n}{2}}\otimes\cdots\otimes|\cdot|^{\frac{n-1}{2}}$ 等于
$\delta^{-\frac12}$。按定义，$\pi(\Delta)$ 有下述简单描述：
\[
\pi(\Delta)=|\cdot|^{\frac{1-n}{2}}\times\cdots\times|\cdot|^{\frac{n-1}{2}}
=\bigl\{\text{光滑函数}\ f\colon B\backslash G\to\mathbb{C}\bigr\}.
\]
因此，唯一的不可约子表示 $Z(\Delta)$ 只能是平凡表示 $\triv$，
它对应于 $B\backslash G$ 上的常值函数。唯一的商 $Q(\Delta)=:\St_n$ 定义为
$\GL_n(K)$ 的（标准）\emph{Steinberg 表示}。它自对偶，即 $\St_n^{\vee}=\St_n$。

\begin{remark}
当 $n=2$、$\Delta=[|\cdot|^{-\frac12},|\cdot|^{\frac12}]$ 时，定理 3.1 的第 (1) 条断言
$\pi(\Delta)$ 长度为 $2$。于是有
\[
0\to\triv\to\pi(\Delta)\to\St_2\to 0 .
\]
\end{remark}

\begin{remark}
另一种做法：考虑 $\Delta_1=[|\cdot|^{\frac12}]$ 与
$\Delta_2=[|\cdot|^{-\frac12}]$，则 $(\Delta_1,\Delta_2)$ 满足 $(*)$ 条件，
且 $\St_2=Z(\Delta_1,\Delta_2)$，$\triv=Q(\Delta_1,\Delta_2)$。作为对比，有短正合列
\[
0\to\triv\to|\cdot|^{-\frac12}\times|\cdot|^{\frac12}\to\St_2\to 0,
\]
\[
0\to\St_2\to|\cdot|^{\frac12}\times|\cdot|^{-\frac12}\to\triv\to 0 .
\]
\end{remark}

\subsubsection{$\GL_2$ 的局部 Langlands 对应}

$n=2$ 的情形格外清楚。在 $\GL_2$ 一侧，$\pi$ 或者超尖点，或者是某个抛物诱导的商
（或者是单个长度为 $2$ 的段，或者是两个长度均为 $1$ 的段）。于是有如下可能：
\begin{enumerate}[label=(\arabic*),leftmargin=2.2em]
\item $\pi$ 是超尖点表示。
\item $\pi=Q(\chi_1,\chi_2)$，其中 $\chi_1/\chi_2\neq|\cdot|^{\pm 1}$，
      于是实际上 $\pi=\chi_1\times\chi_2$——这些称为\emph{主序列表示}。
\item $\pi=Q(\chi|\cdot|^{\frac12},\chi|\cdot|^{-\frac12})$，这是一维表示
      $\triv\otimes\chi$。
\item $\pi=Q(\Delta)$，其中 $\Delta=[\chi|\cdot|^{-\frac12},\chi|\cdot|^{\frac12}]$，
      它等于 $\St_2(\chi):=\St_2\otimes\chi$。
\end{enumerate}
这些称为用 $\chi$ 扭转的 Steinberg 表示（底本作 “Stenberg”，系笔误）。

它们对应的 Weil--Deligne 表示为
\begin{enumerate}[label=(\arabic*),leftmargin=2.2em]
\item 若 $\pi$ 超尖点，则 $\rec_2(\pi)$ 不可约，故
\[
\rec_2(\pi)=(\rho,0),\qquad \rho\ \text{不可约}.
\]
\item 若 $\pi=\chi_1\times\chi_2$ 是主序列、$\chi_1/\chi_2\neq|\cdot|^{\pm 1}$，则
\[
\rec_2(\pi)=(\chi_1\oplus\chi_2,\,0).\en{10}
\]
\item 若 $\pi=\triv\otimes\chi$，则
\[
\rec_2(\pi)=\bigl(\chi\oplus\chi(1),\,0\bigr),
\]
其中 $\chi(1)$ 是 $\chi$ 的「Tate 扭转」。
\item 若 $\pi=\St_2(\chi)$，则
\[
\rec_2(\pi)=\Sp_2(\chi)=\bigl(\chi\oplus\chi(1),\,N\bigr),
\]
即 $N$ 把 $\chi$ 送到 $\chi(1)$，如 1.1.4 的注 (3) 所述。
\end{enumerate}

\subsubsection{把 LLC 归结到超尖点情形的配方}

为证明局部 Langlands 对应（定理 2.13），只需证明存在唯一的一族双射
\[
\rec_n\colon\Asc_n\longrightarrow\Girr_n
\]
满足 2.13 中的全部条件，其中 $\Asc_n$ 是 $\GL_n(K)$ 超尖点表示的等构类（偏序）集，
$\Girr_n$ 是 $W_K$ 的 $n$ 维不可约 Weil--Deligne 表示的等构类集。

这是因为：一旦对所有 $m$ 在 $\Asc_m$ 上有了 $\rec$，对 $\GL_n$ 的不可约容许表示
$\pi=Q(\Delta_1,\dots,\Delta_r)$（Zelevinsky 段 $\Delta_i=\Delta_i(\pi_i,m_i)$，
$n=\sum n_im_i$），就可以定义
\[
\rec_n(\pi)=\bigoplus_{i=1}^{r}\Sp_{m_i}\bigl(\rec_{n_i}(\pi_i)\bigr).
\]
注意 $\rec_n(\pi)=(\rho,N)$ 的底层表示 $\rho$ 正是直和
\[
\rho=\bigoplus_{1\le i\le r}\ \bigoplus_{j=0}^{m_i-1}\rec_{n_i}(\pi_i)(j).
\]
可以直接验证这一构造满足 2.13 的全部要求 (1)--(5)。

\subsection{平方可积、缓增与一般表示}

我们引入如下表示（的包含链）：
\[
\text{超尖点}\subset\text{平方可积}\subset\text{缓增}\subset\text{一般（generic）}
\subset\text{容许}.
\]

\subsubsection{本质 $L^2$ 与本质 $L^{2+\epsilon}$}

\begin{definition}
设 $G^{0}=\{g\in G=\GL_n(K)\mid |\det(g)|_K=1\}$。
\begin{enumerate}[label=(\arabic*),leftmargin=2.2em]
\item 对任意 $r\in\mathbb{R}_{>0}$，容许表示 $(\pi,V)$ 称为\emph{本质 $L^r$}，
      若对任意 $v\in V$、$\lambda\in V^{\vee}$，矩阵系数 $c_{v,\lambda}$ 在 $G^{0}$ 上属于 $L^r$，即
\[
\int_{G^{0}}|c_{v,\lambda}|^{r}\,dg
\]
存在。
\item $(\pi,V)\in\Repadm_G$ 称为 $L^r$，若它本质 $L^r$ 且中心特征标 $\omega_{\pi}$ 酉。
\item $(\pi,V)\in\Repadm_G$ 称为\emph{本质平方可积}（相应地，\emph{平方可积}），
      若它本质 $L^2$（相应地，$L^2$）。
\end{enumerate}
\end{definition}

\begin{remark}
对 $(\pi,V)\in\Repadm_G$，若 $\pi$ 本质 $L^r$，则它对所有 $r'\ge r$ 也都本质 $L^{r'}$。
\end{remark}

\begin{remark}[Casselman 1980]
对任一中心特征标为 $\omega_{\pi}$ 的光滑表示 $(\pi,V)\in\Repinf_G$，
存在（唯一的？）正实值拟特征标 $\chi$，使 $\omega_{\chi\pi}$ 酉。
这样我们总可以通过扭转，把本质 $L^r$ 表示变成 $L^r$ 表示。
\end{remark}

\begin{theorem}[Zelevinsky]
\leavevmode
\begin{enumerate}[label=(\arabic*),leftmargin=2.2em]
\item $\GL_n(K)$ 的不可约容许表示 $\pi$ 本质平方可积，当且仅当它是（单个）Zelevinsky 段
      的 $Q(\Delta)$。
\item 上述本质 $L^2$ 表示 $Q(\Delta)$ 平方可积，当且仅当 $\rho\bigl(\frac{m-1}{2}\bigr)$
      的中心特征标酉，其中 $\Delta=\Delta(\sigma,m)$ 长度为 $m$。
\end{enumerate}
\end{theorem}

\subsubsection{本质缓增与缓增}

我们需要一个比（本质）平方可积稍弱的概念。

\begin{definition}
表示称为\emph{本质缓增}，若对所有 $\epsilon>0$ 它都是本质 $L^{2+\epsilon}$；
若进一步其中心特征标酉，则称为\emph{缓增}。
\end{definition}

\begin{theorem}[Zelevinsky]
$\GL_n(K)$ 的不可约容许表示 $\pi=Q(\Delta_1,\dots,\Delta_r)$ 缓增，
当且仅当每个 $Q(\Delta_i)$ 平方可积。
\end{theorem}

\begin{corollary}
若 $\pi=Q(\Delta_1,\dots,\Delta_r)$ 缓增，则任两个区间 $\Delta_i$ 与 $\Delta_j$ 都不相链；
特别地
\[
\pi\cong Q(\Delta_1)\times\cdots\times Q(\Delta_r).
\]
\end{corollary}

\begin{proof}
设 $\Delta_i=\Delta(\sigma_i,m_i)$。因 $\pi=Q(\Delta_1,\dots,\Delta_r)$ 缓增，
$\sigma_i\bigl(\frac{m_i-1}{2}\bigr)$ 的特征标酉。而若 $\sigma$ 的特征标酉，
则对任一非零 $s\in\mathbb{R}$，$\omega_{\sigma(s)}$ 不酉。
\end{proof}

\begin{remark}
由上述推论，缓增表示恰是那些每个 $Q(\Delta_i)$ 的中心特征标酉的一般表示。
用 Zelevinsky 段来分类本质缓增表示似乎（？）是困难的。
\end{remark}

\subsubsection{Whittaker 模型与一般表示\en{11}}

设 $\psi\colon K\to\mathbb{C}^{\times}$ 是 $K$ 上的非平凡加法拟特征标（如 1.1.10 小节），
$n(\psi)$ 是使 $\psi(\varpi^{-n}\mathcal{O}_K)=1$ 的最大整数（同前）。
设 $U_n\subset\GL_n(K)$ 是严格上三角矩阵构成的子群。定义 $U_n$ 的一维表示 $\theta_{\psi}$：
\[
\theta_{\psi}(u_{ij})=\psi(u_{12}+\cdots+u_{n-1,n}).
\]

\begin{definition}
$\GL_n(K)$ 的不可约容许表示 $\pi$ 称为\emph{一般}（generic），若
$\Hom_{U_n}(\pi|_{U_n},\theta_{\psi})\neq 0$。
\end{definition}

\begin{definition}
设 $(\pi,V)$ 是 $G=\GL_n(K)$ 的不可约容许一般表示，$\pi$ 关于 $\psi$ 的
\emph{Whittaker 模型}定义如下：取非零元素
$\lambda\in\Hom_{U_n}(\pi|_{U_n},\theta_{\psi})$，令
\[
W_{\psi}:=\{f\colon G\to\mathbb{C}\mid f(ug)=\theta_{\psi}(u)f(g)\ 
\ \forall g\in G,\ u\in U_n\}.
\]
$G$ 通过右平移作用在 $W_{\psi}$ 上。定义（单射的）$G$-同态
\[
V\longrightarrow W_{\psi},\qquad v\longmapsto\bigl(g\mapsto\lambda(\pi(g)v)\bigr).
\]
其像 $W_{\pi,\psi}$（同构于 $(\pi,V)$）称为 $\pi$ 的 Whittaker 模型。
\end{definition}

\begin{lemma}
\leavevmode
\begin{enumerate}[label=(\arabic*),leftmargin=2.2em]
\item 「一般」的定义与 $\psi$ 的选取无关。
\item $\pi$ 一般，当且仅当 $\pi^{\vee}$ 一般。
\item 若 $\chi\colon K^{\times}\to\mathbb{C}^{\times}$ 是任一拟特征标，则 $\pi$ 一般
      当且仅当 $\chi\pi$ 一般。
\end{enumerate}
\end{lemma}

\begin{theorem}[Zelevinsky]
$\GL_n(K)$ 的不可约容许表示 $\pi=Q(\Delta_1,\dots,\Delta_r)$ 一般，
当且仅当任两个 Zelevinsky 段 $\Delta_i$ 与 $\Delta_j$ 都不相链；由定理 3.1 第 (5) 条，
这又等价于
\[
\pi=Q(\Delta_1)\times\cdots\times Q(\Delta_r).
\]
特别地，推论 3.4 可以重新表述为：每个不可约（本质）缓增表示 $(\pi,V)$ 都是一般的。
\end{theorem}

\begin{remark}[Shalika]
一般表示大体上就是整体应用中我们所关心的表示。设 $F$ 是数域，则任一不可约容许
尖点表示 $\pi=\otimes'_v\pi_v$ 的所有局部分量 $\pi_v$ 都是一般的。
\end{remark}

\subsection{$K$-球面表示}

设 $K\subset G=\GL_n(K)$ 是紧开子群，不可约容许表示 $\pi$ 称为 \emph{$K$-球面}的，
若 $\pi^{K}\neq 0$；若对极大紧开子群 $K_0=\GL_n(\mathcal{O}_K)$ 是 $K_0$-球面的，
则称 $\pi$ 为\emph{球面的}。本小节考虑球面表示与 Iwahori-球面表示。

\subsubsection{$\GL_n(K)$ 的结构}

记号：
\begin{itemize}[leftmargin=1.6em]
\item 设 $B\subset G=\GL_n(K)$ 是标准 Borel 子群，$U\subset B$ 是严格上三角矩阵，
      $U^{-}$ 是严格下三角矩阵。设 $T\subset B$ 是对角环面（$B$ 相应的 Levi）。
\item $K_0=\GL_n(\mathcal{O}_K)$ 是极大紧子群。
\item 设 $T_0:=T\cap K_0\cong(\mathcal{O}_K^{\times})^{n}$。
\item 设 $B_0\subset K_0$ 是 Iwahori 子群，即标准 Borel $B(k)\subset\GL_n(k)$
      在约化映射下的原像。在本文中我们交替使用 Iw $=B_0$。
\item 设 $N=N_G(T)\subset G$ 是（广义）置换矩阵群，$N_K=N\cap K$。
      我们把 $N_K/T_0$ 与 Weyl 群 $W$ 等同。
\item 设 $W_1:=N/T_0$，有时称为\emph{广义 Weyl 群}。
\end{itemize}

回忆下述分解。

\begin{lemma}
\leavevmode
\begin{enumerate}[label=(\arabic*),leftmargin=2.2em]
\item $\displaystyle K_0=\bigsqcup_{w\in W}B_0wB_0$。
\item $G$ 的 Iwasawa 分解：$G=BK_0$。
\item Cartan 分解：设 $T^{-}\subset T$ 是对角矩阵
      $\mathrm{Diag}(\varpi_K^{r_1},\varpi_K^{r_2},\dots,\varpi_K^{r_n})$（$r_1\ge r_2\ge\cdots\ge r_n$）
      构成的子集，则
\[
G=\bigsqcup_{\lambda\in T^{-}/T^{0}}K_0\lambda K_0 .
\]
\item Bruhat--Tits 分解：仍记 $W_1=N/T_0$，则
\[
G=\bigsqcup_{w_1\in W_1}B_0w_1B_0 .
\]
\item Iwahori 分解：$B_0=(B_0\cap U^{-})\cdot T_0\cdot(B_0\cap U)$。
\end{enumerate}
\end{lemma}

讨论 Hecke 算子时，把 $K_0\lambda K_0$ 写成一族互不相交的 $\lambda'K_0$ 之并
也有用；初稿中我们跳过这一点。

\subsubsection{非分歧主序列}

底本此处为空（仅标题，作者原拟在二稿补写）。

\subsubsection{非分歧表示}

$G=\GL_n(K)$ 的每个不可约非分歧表示 $\pi$ 都形如 $Q(\chi_1,\dots,\chi_n)$，
其中 $\chi_1,\dots,\chi_n$ 是 $K^{\times}$ 的非分歧拟特征标。

\begin{example}
考虑 $\GL_2(K)$ 的情形。上述断言说明 $\GL_2(K)$ 的不可约非分歧表示是下列之一：
\begin{enumerate}[label=(\arabic*),leftmargin=2.2em]
\item $\pi=\chi_1\times\chi_2$，其中 $\chi_1/\chi_2\neq|\cdot|^{\pm 1}$，
      是 $\chi_1,\chi_2$ 为非分歧拟特征标的主序列表示；
\item $\pi=Q(\chi|\cdot|^{\frac12},\chi|\cdot|^{-\frac12})=\triv\otimes\chi$，
      其中 $\chi$ 是非分歧拟特征标。
\end{enumerate}
\end{example}

\subsubsection{Iwahori-不动向量}

底本此处写作「to be added in second draft」（留待二稿补写）。

\subsection{容许表示的导子与新向量}

底本此节只有标题，正文未写。

\subsection{深度 $0$ 的表示}

我们把这一讨论总结为下图。这里只关心自守一侧的不可约容许表示，以及
Weil--Deligne 表示 $\rec_n(\pi)=(\rho,N)$（或相应的 $l$-进表示 $r_n(\pi)$，
使 $\rec_n(\pi)=\mathrm{WD}(r_n(\pi))$）。

\begin{center}
\renewcommand{\arraystretch}{1.35}
\begin{tabular}{@{}r@{\qquad$\Longleftrightarrow$\qquad}l@{}}
$\pi$ 非分歧 & $\pi^{K}\neq 0$\\
$r_n(\pi)$ 非分歧 & $\rho|_{I_K}=0$ 且 $N=0$\\
$\pi$ 是 Iwahori-球面的 & $\pi^{\mathrm{Iw}}\neq 0$\\
$\rec_n(\pi)$ 为 $I_K$-球面的 & $\rho|_{I_K}=0$\\
$\pi$ 的深度为 $0$ & $r_n(\pi)$ 驯分歧
\end{tabular}
\end{center}

\subsubsection{导子与深度}

底本此处只有标题，正文未写。

\subsection{$L$-因子与 $\epsilon$-因子}

底本此处写作「To be completed in second draft」（留待二稿补完）。

\subsubsection{一般表示的 $\epsilon$-因子}

我们把这一讨论留到将来的草稿，但先记下如下几点。
表示论一侧的 $\epsilon$-因子，是以超尖点情形的 $\epsilon$-因子为基础、
依托 Bernstein--Zelevinsky 分类来定义的。对一对表示 $\pi,\pi'$，我们将定义
$\epsilon(\pi\times\pi',s,\psi)$。

\subsubsection{归纳定义}

底本此处只有标题，正文未写。

\subsubsection{$\epsilon$-因子与导子}

设 $\pi$ 是 $\GL_n(K)$ 的光滑不可约表示。当 $\pi'\colon K^{\times}\to\mathbb{C}^{\times}$
取平凡特征标时，我们用 $L(\pi,s)$ 与 $\epsilon(\pi,s,\psi)$ 记相应的 $L$-因子与
$\epsilon$-因子。仍记 $q=\#k$。由 JPPS 的定理（及其勘误），有
\[
\epsilon(\pi,s,\psi)=q^{-\left(c(\pi)+n\cdot n(\psi)\right)\cdot s}\cdot
\epsilon(\pi,0,\psi),
\]
其中 $c(\pi)$ 是 3.4 小节定义的 $\pi$ 的导子。（特别地，文献中 $\pi$ 的导子的
两个定义——其一是 3.4 小节的定义，其二是由 $L(\pi,s)$ 的函数方程给出的定义——
彼此一致。）

\subsubsection{导子被 LLC 保持}

现在有下述平凡的观察。

\begin{lemma}
在局部 Langlands 对应之下，有
\[
c(\pi)=c\bigl(\rec_n(\pi)\bigr).
\]
换言之，导子被保持。
\end{lemma}

\begin{proof}
取 $\pi'$ 为平凡特征标，由 $\epsilon$-因子的相容性得
$\epsilon(\pi,\psi,s)=\epsilon(\rec_n(\pi),\psi,s)$，特别地，由推论 1.4 及上述讨论，
\[
q^{-\left(c(\pi)+n\cdot n(\psi)\right)\cdot s}
=q^{-\left(c(\rec_n(\pi))+n\cdot n(\psi)\right)\cdot s}.
\]
\end{proof}

\begin{keytakeaways}
\begin{itemize}[leftmargin=1.5em]
\item Bernstein--Zelevinsky 分类：$\GL_n$ 的不可约容许表示唯一地写成
      $Q(\Delta_1,\dots,\Delta_r)$，段族在 $(*)$-置换下唯一；$\pi(\Delta)$ 的合成列长度 $2^{m-1}$。
\item 该分类给出 $R_K=\bigoplus_n R_n(K)$ 的分次交换环结构与对合 $Q\mapsto Z$，
      并给出基本构件：超尖点（单段且段长 $1$ 时）、本质平方可积（单段）、
      缓增（各段平方可积）、一般（各段互不相链）。
\item LLC 的证明可以归结为超尖点情形：$\rec_n(\pi)=\bigoplus_i\Sp_{m_i}(\rec_{n_i}(\pi_i))$，
      其余性质自动成立；这就是 Henniart 论证的骨架（3.1.7 与 2.3.1 合用）。
\item $\GL_2$ 的四类表示（超尖点、主序列、一维、Steinberg 扭转）分别对应
      $(\rho,0)$ 不可约、$(\chi_1\oplus\chi_2,0)$、$(\chi\oplus\chi(1),0)$、$\Sp_2(\chi)$。
\item 导子被 LLC 保持：$c(\pi)=c(\rec_n(\pi))$，由 $\epsilon$-因子的相容性与
      JPPS 公式 $q^{-(c(\pi)+n\,n(\psi))s}$ 直接得到。
\end{itemize}
\end{keytakeaways}

% 4. 惯性 Langlands 对应（底本 p.32--35）
\section{惯性 Langlands 对应}

\subsection{Bernstein 分解}

\subsubsection{Bernstein 分量与分解}

\begin{definition}
设 $L\subset\GL_n(K)$ 是 Levi 子群，$P\subset\GL_n(K)$ 是以 $L$ 为 Levi 因子的抛物子群。
设 $\varsigma$ 是 $L$ 的不可约超尖点表示。\emph{Bernstein 分量}
\[
R_{(L,\varsigma)}=R_{(L,\varsigma)}\bigl(\GL_n(K)\bigr)\subset\Repinf_n
\]
是这样一个满子范畴：$\pi$ 属于它，当且仅当 $\pi$ 的每个不可约子商都是
\[
I^{\GL_n}_L(\varsigma\otimes\eta):=\Ind^{\GL_n(K)}_P(\varsigma\otimes\eta)
\]
的子商，其中 $\eta$ 是 $L$ 的非分歧特征标，而 $I^{\GL_n}_L:=\Ind^{\GL_n(K)}_P\circ j^{*}$
是正规化抛物诱导，$j\colon P\to L$ 是到 Levi 商的映射。
\end{definition}

\begin{lemma}
两个 Bernstein 分量 $R_{(L_1,\varsigma_1)}$ 与 $R_{(L_2,\varsigma_2)}$ 相等，
当且仅当它们\emph{惯性等价}，即存在 $L_2$ 的非分歧特征标 $\eta$ 与 $g\in\GL_n(K)$ 使
\[
gL_1g^{-1}=L_2,\qquad g\varsigma_1=\varsigma_2\eta .
\]
\end{lemma}

\begin{definition}
定义 $\GL_n(K)$ 的\emph{Bernstein 中心}为 Bernstein 分量的集合，即
\[
\mathcal{B}=\mathcal{B}\bigl(\GL_n(K)\bigr):=\{(L,\varsigma)\}/\text{惯性等价}.
\]
\end{definition}

\begin{theorem}[Bernstein 分解]
$\GL_n(K)$ 的光滑表示范畴 $\Repinf_n$ 分解为直积
\[
\Repinf_n=\prod_{(L,\varsigma)\in\mathcal{B}}R_{(L,\varsigma)} .
\]
\end{theorem}

\begin{remark}
要确定一个给定的不可约表示 $\pi\in\An(K)$ 落在哪个 Bernstein 分量中，
只需考虑它的超尖点支撑。
\end{remark}

\subsubsection{$\GL_2$ 的 Bernstein 分量}

\begin{example}[$n=2$]
Bernstein 分量 $R_{(L,\sigma)}(\GL_2(K))$ 是下列类型之一：
\begin{enumerate}[label=(\arabic*),leftmargin=2.2em]
\item 若 $R_{(L,\sigma)}(G)$ 含有不可约超尖点表示 $\pi$，则 $(L,\sigma)$ 与 $(G,\pi)$
      惯性等价，而 Bernstein 分量 $R_{(G,\pi)}(G)$ 中的不可约表示为
\[
\{\pi\otimes(\eta\circ\det):\ \eta\ \text{非分歧}\}.
\]
\item 若 $R_{(L,\sigma)}(G)$ 含有 $\pi=\St(\chi)=\St_G\otimes(\chi\circ\det)$，
      则 $(L,\sigma)\sim(T,\chi\otimes\chi)$，而 $R_{(T,\chi)}(G)$ 中的不可约表示为
\[
\bigl\{\St_G\otimes(\chi\eta\circ\det)\bigr\}
\ \cup\
\bigl\{(\eta_1\chi|\cdot|^{-1/2})\times(\eta_2\chi|\cdot|^{1/2})\bigr\},
\]
其中 $\eta,\eta_1,\eta_2$ 取遍非分歧特征标（且使后一族仍为主序列）。
\item 若 $\pi=\pi(\chi_1,\chi_2)$，其中 $\chi_1\chi_2^{-1}$ 分歧
      （否则落入前一情形），则该 Bernstein 分量中的不可约表示为
\[
\{(\eta_1\chi_1)\times(\eta_2\chi_2)\},
\]
其中 $\eta_1,\eta_2$ 非分歧。
\end{enumerate}
\end{example}

\begin{remark}
在差一个非分歧特征标的意义下，上例中各个 Bernstein 分量对应的超尖点支撑为：
\begin{enumerate}[label=(\arabic*),leftmargin=2.2em]
\item $R_{(G,\pi)}(G)$ 对应 $\{\pi\}$；
\item $R_{(T,\chi\otimes\chi)}(G)$ 对应 $\{\chi,\chi\}$；
\item $R_{(T,\chi_1\otimes\chi_2)}(G)$ 对应 $\{\chi_1,\chi_2\}$，其中 $\chi_1\chi_2^{-1}$ 分歧。
\end{enumerate}
\end{remark}

\subsubsection{光滑表示的超尖点部分}

底本此处只有标题，正文未写。

\subsubsection{Bernstein 分解的证明}

底本此处只有标题，正文未写。

\subsection{惯性对应}

\subsubsection{Weil--Deligne 惯性类型}

先需要一个（有限维）表示论中的标准引理：

\begin{lemma}
设 $H\triangleleft G$ 是正规子群，$G/H$ 交换，$(\rho_1,V_1)$ 与 $(\rho_2,V_2)$ 是
$\mathbb{C}$ 上（或任一代数闭域上）两个不可约有限维表示，且 $\rho_1|_H\cong\rho_2|_H$，
则存在特征标 $\chi\colon G\to G/H\to\mathbb{C}^{\times}$，使 $\rho_2\cong\rho_1\otimes\chi$。
\end{lemma}

\begin{proof}[略证]
把 $\Hom_{\mathbb{C}}(V_1,V_2)$ 看作带有 $G$-作用的向量空间。由
$\rho_1|_H\cong\rho_2|_H$ 知
\[
W:=\Hom_H(\rho_1,\rho_2)=\bigl(\Hom_{\mathbb{C}}(V_1,V_2)\bigr)^{H}\neq 0 .
\]
$W\subset\Hom_{\mathbb{C}}(V_1,V_2)$ 是 $G$ 的子表示，且 $H$ 在其上平凡作用，
于是交换商 $G/H$ 作用在 $W$ 上；$W$ 因此含有一个一维 $G$-子表示：
\[
\chi\hookrightarrow W\subset V_1^{\vee}\otimes V_2 .
\]
特别地，用 $\chi$ 扭转可得
\[
\triv\hookrightarrow(V_1\otimes\chi)^{\vee}\otimes V_2,
\]
即
\[
\triv\hookrightarrow\bigl((V_1\otimes\chi)^{\vee}\otimes V_2\bigr)^{G}
=\Hom_G(V_1\otimes\chi,V_2)\neq 0 .
\]
于是由 Schur 引理，$\rho_1\chi\cong\rho_2$。
\end{proof}

\begin{remark}
事实上我们证明了稍强一些的断言，下文会用到：在上述设定中，若 $\rho_1$ 与 $\rho_2$
都不可约，且 $\Hom_H(\rho_1,\rho_2)\neq 0$，则同样的结论成立。
也就是说，引理中并不需要事先假设 $\rho_1|_H\cong\rho_2|_H$。
\end{remark}

\begin{lemma}
$G$ 的两个不可约容许表示 $\pi$ 与 $\pi'$ 属于同一个 Bernstein 分量
$R_{(L,\sigma)}(G)$，当且仅当由 LLC 得到的（Weil--Deligne 表示的）底层 $W_K$-表示满足
\[
\rec_n(\pi)|_{I_K}\cong\rec_n(\pi')|_{I_K}.
\]
换言之，Bernstein 中心 $\mathcal{B}(\GL_n)$ 由 Weil--Deligne 惯性类型标定。
\end{lemma}

\begin{proof}
由 Bernstein--Zelevinsky 分类，$\pi=Q(\Delta_1,\dots,\Delta_r)$，
其中 $\Delta_i=\Delta_i(\pi_i,m_i)$ 是 Zelevinsky 段，于是 $\pi$ 属于这样的 Bernstein 分量：
\[
L=\prod_{i=1}^{r}\GL_{n_i}(K)^{m_i},
\qquad
\sigma=\bigotimes_{i=1}^{r}\ \bigotimes_{j=0}^{m_i-1}\pi_{i,j},
\]
其中 $\pi_{i,j}$ 是 $\GL_{n_i}(K)$ 的超尖点表示 $\pi_i|\cdot|^{j}$。由局部 Langlands 对应，
\[
\rec_n(\pi)=\bigoplus_{i=1}^{r}\Sp_{m_i}\bigl(\rec_{n_i}(\pi_i)\bigr),
\]
而限制到 $I_K$ 的底层 $W_K$-表示为
\[
\rec_n(\pi)|_{I_K}=\bigoplus_{i=1}^{r}\rec_{n_i}(\pi_i)_{m_i}|_{I_K}.
\]
于是归结到超尖点情形：$\GL_k(K)$ 的两个超尖点表示 $\pi$ 与 $\pi'$ 属于同一
Bernstein 分量（即二者相差一个非分歧特征标 $\eta$ 的扭转），当且仅当
$\rec(\pi)|_{I_K}\cong\rec(\pi')|_{I_K}$，而这可由引理 4.3 推出。
注意最后一步隐含地用到了引理 4.3 之后的注：设
\[
\bigoplus_{i=1}^{r}\rec_{n_i}(\pi_i)_{m_i}|_{I_K}\cong
\bigoplus_{i=1}^{r'}\rec_{n'_i}(\pi'_i)_{m'_i}|_{I_K},
\]
我们（逐次）应用该注，可得对某些非分歧特征标 $\psi_i$ 有
$\rec_{n_i}(\pi_i)\otimes\psi_i\cong\rec_{n_i}(\pi'_i)$。
\end{proof}

\subsubsection{Galois 惯性类型与一般表示}

设 $\mathcal{K}=\GL_n(\mathcal{O}_K)$ 为 $\GL_n(K)$ 的标准极大紧开子群。
（底本此处仍写作 $K=\GL_n(\mathcal{O}_K)$，与域 $K$ 撞记号；译文改用 $\mathcal{K}$。）

\begin{lemma}
设 $\pi,\pi'$ 是 $\GL_n(K)$ 的两个一般不可约容许表示，其相应的 $l$-进表示满足
\[
r_n(\pi)|_{I_K}\cong r_n(\pi')|_{I_K},
\qquad r_n=\mathrm{WD}^{-1}\circ\rec_n .
\]
则
\[
\pi|_{\mathcal{K}}\cong\pi'|_{\mathcal{K}} .
\]
换言之，若两个一般表示有相同的 Galois 惯性类型，则它们限制到 $\mathcal{K}$ 上同构。
\end{lemma}

\begin{proof}
由 Iwasawa 分解 $\GL_n(K)=B\cdot\GL_n(\mathcal{O}_K)$（$B$ 为上三角矩阵构成的 Borel 子群）
以及诱导表示限制的 Mackey 分解，对任一（标准）抛物子群 $P\subset\GL_n(K)$
与 $P$ 的表示 $\sigma$，有
\[
\bigl(\Ind^{\GL_n(K)}_P\sigma\bigr)|_{\mathcal{K}}
=\Ind^{\mathcal{K}}_{P\cap\mathcal{K}}\bigl(\sigma|_{P\cap\mathcal{K}}\bigr).
\]
因 $\pi$ 一般，由定理 3.1 有 $\pi=Q(\Delta_1)\times\cdots\times Q(\Delta_r)$，且
$\rec_n(\pi)=\bigoplus_i\rec_{n_im_i}(Q(\Delta_i))$，$\pi'$ 同理。$W_K$ 的两个表示
\[
\bigoplus_{i=1}^{r}\Sp_{m_i}(\rho_i),
\qquad
\bigoplus_{i=1}^{r'}\Sp_{m'_i}(\rho'_i)
\]
有相同的 Galois 惯性类型，当且仅当由引理 1.9，$r=r'$，且在重排之后
$m_i=m'_i$、$\rho_i|_{I_K}\cong\rho'_i|_{I_K}$。

于是只需在 $\pi=Q(\Delta(\lambda,m))$ 与 $\pi'=Q(\Delta(\lambda',m))$ 的假设下证明本引理，
其中 $\lambda,\lambda'$ 是 $\GL_d(K)$ 的超尖点表示。注意 $\lambda=\lambda'\otimes\eta$，
其中 $\eta$ 是 $\GL_d(K)$ 的非分歧特征标。特征标 $\eta$ 形如 $\eta=\chi\circ\det$，
$\chi$ 是 $K^{\times}\to\mathbb{C}^{\times}$ 的非分歧拟特征标。于是
$\rec_n(\pi)=\rec_n(\pi')\otimes\rec_1(\chi)$，由此得 $\pi=\pi'\chi$。
\end{proof}

\begin{remark}
「一般」这一条件不可省。例如对 $\GL_2(K)$，$\triv$ 与 $\St_2$ 落在同一
Bernstein 分量中，但它们限制到 $\mathcal{K}=\GL_2(\mathcal{O}_K)$ 上并不同构。
\end{remark}

\subsection{类型}

最后我们简述类型论。

\subsubsection{标准定义}

\begin{definition}
设 $\sigma$ 是 $\mathcal{K}$ 的光滑不可约表示，$s=(L,\varsigma)\in\mathcal{B}(\GL_n)$
是 Bernstein 中心中的一个元素。
\begin{enumerate}[label=(\arabic*),leftmargin=2.2em]
\item 若对一切光滑不可约表示 $\pi\in\Repinf_n$，
\[
\pi|_{\mathcal{K}}\ \text{包含}\ \sigma
\quad\Longleftrightarrow\quad
\pi\in R_{(L,\varsigma)},
\]
则称 $\sigma$ 是 $(L,\varsigma)$ 的\emph{类型}（type）。
\item 若对一切光滑不可约表示 $\pi$，
\[
\pi|_{\mathcal{K}}\ \text{包含}\ \sigma
\quad\Longrightarrow\quad
\pi\in R_{(L,\varsigma)},
\]
则称 $\sigma$ 是 $(L,\varsigma)$ 的 \emph{$\mathcal{K}$-型}。
\item 设 $\pi\in R_{(L,\varsigma)}$ 是 $R_{(L,\varsigma)}$ 中的光滑不可约表示，
      若 $\sigma$ 是 $(L,\varsigma)$ 的 $\mathcal{K}$-型，且在 $\pi|_{\mathcal{K}}$ 中出现的
      $\mathcal{K}$-型里维数最小，则称 $\sigma$ 是 $\pi$ 的\emph{极小 $\mathcal{K}$-型}。
\end{enumerate}
\end{definition}

注意每个类型显然都是 $\mathcal{K}$-型，且这两个概念只依赖于
$(L,\varsigma)\in\mathcal{B}(\GL_n)$；而 $\pi$ 的极小 $\mathcal{K}$-型的定义依赖于
特定的 $\pi\in R_{(L,\varsigma)}$。也就是说，对给定的 $s=(L,\varsigma)$ 可以有很多
$\mathcal{K}$-型；对每个不可约 $\pi\in R_{(L,\varsigma)}$，一个 $\mathcal{K}$-型 $\sigma$
可能在 $\pi|_{\mathcal{K}}$ 中出现、也可能不出现；在那些出现的当中，维数最小者
就是 $\pi$ 的极小 $\mathcal{K}$-型。

\begin{remark}
我们对「类型」的定义与 [9] 中的定义一致，而与 [4] 中的定义略有不同。
后者我们称为 \emph{Bushnell--Kutzko 型}，它是一个对 $(J,\sigma)$，
其中 $J$ 是紧开子群、$\sigma$ 是 $J$ 的光滑不可约表示——
也就是说，它允许比 $\mathcal{K}=\GL_n(\mathcal{O}_K)$ 更一般的紧开子群。
\end{remark}

\begin{keytakeaways}
\begin{itemize}[leftmargin=1.5em]
\item Bernstein 分解把 $\Repinf_n$ 拆成直积 $\prod_{(L,\varsigma)}R_{(L,\varsigma)}$，
      分量的标定是「惯性等价」类 $(L,\varsigma)$，即超尖点支撑模去非分歧扭转。
\item 惯性 Langlands 对应（引理 4.4）：$\pi$ 与 $\pi'$ 在同一 Bernstein 分量
      $\Leftrightarrow$ $\rec_n(\pi)|_{I_K}\cong\rec_n(\pi')|_{I_K}$——
      即 Bernstein 中心由 Weil--Deligne 惯性类型标定。
\item 一般表示更进一步（引理 4.5）：Galois 惯性类型相同 $\Rightarrow$
      $\pi|_{\mathcal{K}}\cong\pi'|_{\mathcal{K}}$；「一般」不可去掉。
\item 类型（type）、$\mathcal{K}$-型与极小 $\mathcal{K}$-型是把 Bernstein 分量
      在 $\GL_n(\mathcal{O}_K)$ 一侧「实现」出来的工具，与 Bushnell--Kutzko 型的
      区别在于是否允许一般紧开子群。
\end{itemize}
\end{keytakeaways}

% 5--6：底本只有标题（p.36）
\section{Carayol 的上同调实现纲领}

底本此节只有标题，正文未写。

\section{局部 Langlands 对应的族}

底本此节只有标题，正文未写。

\begin{keytakeaways}
\begin{itemize}[leftmargin=1.5em]
\item 这两节在底本中仅存标题，说明作者原拟讨论的两条后续路线：
      一是 Carayol 式的「用 Shimura 簇（或相关代数簇）的中部上同调实现 LLC」，
      二是「把 LLC 组织成参数空间上的族」（即后来的族与叠的视角）。
\item 本译注本按体例照录标题、不代拟内容；读者若需相关内容，
      可参考本文参考文献 [8]（Harris--Taylor）与 [11]（Wedhorn 讲义）。
\end{itemize}
\end{keytakeaways}

% 参考文献与底本原注（p.36）
\section*{参考文献}
\addcontentsline{toc}{section}{参考文献}

\begin{enumerate}[label={[\arabic*]},leftmargin=3.2em,itemsep=0.15em]
\item J. Bernstein. \emph{Harvard lecture notes on representations of $p$-adic reductive groups.}
      （$p$ 进约化群表示论的哈佛讲义。）
\item J. Bernstein, P. Deligne. \emph{Le centre de Bernstein.}
      In: \emph{Représentations des groupes réductifs sur un corps local},
      Travaux en cours (P. Deligne ed.), Hermann, Paris, 1--32 (1984).
      （Bernstein 中心。）
\item J. Bernstein, A. V. Zelevinsky. \emph{Induced representations of reductive $p$-adic groups. I.}
      Ann. Sci. École Norm. Sup. \textbf{10}, 441--472 (1977).
      （约化 $p$ 进群的诱导表示 I。）
\item Bushnell, Kutzko. \emph{The Admissible Dual of $GL(N)$ via Compact Open Subgroups.}
      Annals of Mathematics Studies, 1993.
      （经由紧开子群给出的 $GL(N)$ 的容许对偶。）
\item P. Cartier. \emph{Representations of $p$-adic groups: a survey.}
      Proc. of Symp. in Pure Math. \textbf{33}, part 1, 111--155 (1979).
      （$p$ 进群表示论综述。）
\item W. Casselman. \emph{The unramified principal series of $p$-adic groups I.}
      Compositio Math. \textbf{40}, no. 3, 387--406 (1980).
      （$p$ 进群的非分歧主序列 I。）
\item G. Henniart. \emph{Sur l'unicité des types pour $GL_2$.}
      Duke Mathematical Journal \textbf{2}, no. 115 (2002). Appendix to
      \emph{Multiplicités modulaires et représentations de $GL_2(\mathbb{Z}_p)$
      et de $\mathrm{Gal}(\mathbb{Q}_p/\mathbb{Q}_p)$ en $l=p$}.
      （论 $GL_2$ 类型的唯一性；附录于「模重数与 $l=p$ 时 $GL_2(\mathbb{Z}_p)$
      及 $\mathrm{Gal}(\overline{\mathbb{Q}}_p/\mathbb{Q}_p)$ 的表示」。）
\item M. Harris, R. Taylor. \emph{The geometry and cohomology of some simple Shimura varieties.}
      Annals of Mathematics Studies, vol. 151 (2001).
      （某些简单 Shimura 簇的几何与上同调。）
\item V. Păsku̧nas. \emph{Unicity of types for supercuspidal representations of $GL_N$.}
      Proc. London Math. Soc. (3) \textbf{91}, no. 3, 623--654 (2005).
      （$GL_N$ 超尖点表示类型的唯一性。）
\item J. Tate. \emph{Number theoretic background.}
      In: \emph{Automorphic forms, representations, and $L$-functions}, vol. 2 (1977).
      （数论背景。）
\item T. Wedhorn. \emph{The Local Langlands Correspondence for $GL(n)$ over $p$-adic Fields.}
      Lecture notes (2000).
      （$p$ 进域上 $GL(n)$ 的局部 Langlands 对应。）
\end{enumerate}

\vspace{0.8em}
\noindent\textbf{作者通信地址：}\ \texttt{zijian.yao.math@gmail.com}

\bigskip
\hrule
\bigskip

\section*{原注}
\addcontentsline{toc}{section}{原注}

底本以页脚脚注形式给出 11 条注，正文以上标数字标出。译文照录其位置与内容如下；
凡底本在版式上被截断之处，如实标出，不作补写。

\begin{enumerate}[label={\textsuperscript{\arabic*}},leftmargin=2.4em,itemsep=0.3em]
\item （p.11，$H(G)$ 的定义处）不要与下一小节所称的 Hecke 作用混淆。
\item （p.13）例如，这正是证明 $H(G)=\bigcup_K H(G/\!/K)$ 的方法；或者等价地说，
      两个左、右[不变子空间]……\emph{（底本此注在页脚处被排版截断，末句未完。）}
\item （p.14）在 $K^{n}$ 上引入非退化对称双线性形式之后。
\item （p.14）注意我们其实不能直接应用定理 2.5，因为它的证明用到了 Schur 引理。
\item （p.15）由于时间限制，我作报告时省略了这部分讨论。
\item （p.17）这在 Cartier 的综述 [5] 中称为「对偶 Frobenius 互反」。
\item （p.18）更精确地说，我们应取 $K$-抛物子群 $P\subset G$，然后取其 $K$-点。
\item （p.20）我认为条件 (2)（等价地 (3)）可能正是超尖点性最初的定义方式。
\item （p.23）文献中有时称之为 Bernstein--Zelevinsky 段；本文中我们只用
      「段」或「Zelevinsky 段」。
\item （p.26）更准确地说是 $(\Art^{-1}\chi_1\oplus\Art^{-1}\chi_2,\,0)$。
\item （p.28）二稿中将补充大量内容，尤其是 Whittaker 模型及其在定义
      $\epsilon$[-因子]中的重要性。\emph{（底本此注在页脚处被排版截断。）}
\end{enumerate}


\end{document}
